\documentclass[12pt]{article}
\usepackage[margin=1in]{geometry}
\usepackage{amsmath,amssymb,amsfonts,amsthm,bm,array,booktabs,enumitem}
\usepackage[colorlinks=true,linkcolor=blue,citecolor=blue,urlcolor=blue]{hyperref}

\numberwithin{equation}{section}

\newtheorem{theorem}{Theorem}
\newtheorem{proposition}{Proposition}
\newtheorem{lemma}{Lemma}
\newtheorem{assumption}{Assumption}
\theoremstyle{definition}
\newtheorem{definition}{Definition}
\newtheorem{remark}{Remark}

\newcommand{\E}{\mathbb E}
\newcommand{\R}{\mathbb R}
\newcommand{\Pp}{\mathbb P}
\newcommand{\ind}{\mathbf 1}
\newcommand{\op}{o_p}
\newcommand{\Op}{O_p}
\newcommand{\N}{\mathcal N}
\newcommand{\LtwoT}{L^2(P_T)}
\newcommand{\LtwoP}{L^2(P)}
\newcommand{\innerT}[2]{\langle #1,#2\rangle_T}
\newcommand{\innerP}[2]{\langle #1,#2\rangle_P}

\begin{document}

\begin{center}
{\Large\bf L2 Geometry Workbook for Profile Sieve EL}\\[6pt]
{\large Main-Body Trial: Geometry-First Proof With Detailed Appendix}\\[6pt]
{\normalsize Trial structure: main proof first, detailed machinery retained later}\\[12pt]
\end{center}

\tableofcontents
\newpage

\part{Main Body Trial: Geometry-First Proof}


\section{Model Setup: Full Geometry in \texorpdfstring{\(L^2(P_T)\)}{L2PT} and \texorpdfstring{\(L^2(P)\)}{L2P}}

The purpose of this workbook is to rewrite the whole proof logic as Hilbert
space geometry. The sample space is not first probability; it is first the
finite empirical Hilbert space
\[
    \LtwoT=(\R^T,\innerT{a}{b}=T^{-1}a'b).
\]
The population limit is the Hilbert space
\[
    \LtwoP,\qquad \innerP{f}{g}=\E\{fg\}.
\]

The model setup is the following chain:
\[
\begin{aligned}
\bm y&=\beta_0\bm x+\bm\mu_0+\bm u
&&\text{model in }\LtwoT,\\
\mu_{0,t}&=m_{j_0(t)}(W_{t-1}),\quad
j_0(t)=j\iff k_{j,0}<t\le k_{j+1,0}
&&\text{true nuisance vector},\\
\N_{K,\Lambda,T}&=\operatorname{col}(\bm Q_{K,\Lambda})
&&\text{nuisance subspace},\\
\bm M_{K,\Lambda}&=I-\Pi_{K,\Lambda,T}
&&\text{residual-maker},\\
S_T(\beta,\Lambda)&=\sqrt T\innerT{\bm M_{K,\Lambda}(\bm y-\beta\bm x)}
{\bm M_{K,\Lambda}\bm w}
&&\text{projected score}.
\end{aligned}
\]
Here \(\bm\mu_0\) is exactly the stacked nonparametric nuisance part of the
model. If there are no nuisance breaks, \(\mu_{0,t}=m(W_{t-1})\). With breaks,
\(\mu_{0,t}=m_j(W_{t-1})\) on regime \(j\), so \(\bm\mu_0\) is not one stable
function evaluated over the whole sample; it is the true piecewise nuisance
surface evaluated observation by observation.

EL then uses the scalar products
\[
    Z_t(\beta,\Lambda)
    =\{\bm M_{K,\Lambda}(\bm y-\beta\bm x)\}_t\{\bm M_{K,\Lambda}\bm w\}_t
\]
to form a one-dimensional empirical likelihood. EL is not the Hilbert
projection; EL is the likelihood wrapper around the projected score.

\begin{center}
\fbox{\parbox{0.92\textwidth}{
\textbf{Main-body proof map.} The front proof states the finite-sample geometry needed for the main theorem and then proves consistency and Wilks on a good event \(\mathcal G_T\). The detailed projection algebra, stochastic certification, and older merge derivations are retained in the appendix-style supporting material. Probability builds \(\mathcal G_T\); geometry carries the proof once \(\mathcal G_T\) holds.}}
\end{center}

\section{Unified Geometric Dictionary}

This workbook uses one geometric dictionary throughout. The partition does not
only index a design matrix. It indexes a nuisance subspace:
\[
    \mathfrak N_T:\Lambda\longmapsto \N_{K,\Lambda,T}
    =\operatorname{col}(\bm Q_{K,\Lambda})\subset\LtwoT.
\]
Thus break search is a finite-dimensional subspace-selection problem. For each
\(\Lambda\), the profiled nuisance estimate is the foot of the perpendicular
from \(\bm y_\beta^\circ=\bm y-\beta\bm x\) to \(\N_{K,\Lambda,T}\), and the
profile residual
\[
    \bm M_{K,\Lambda}\bm y_\beta^\circ
\]
lives in the normal space
\[
    \N_{K,\Lambda,T}^{\perp_T}.
\]
The raw score direction \(\bm w\) is also pushed into the same normal space:
\[
    \bm M_{K,\Lambda}\bm w\in \N_{K,\Lambda,T}^{\perp_T}.
\]
Therefore inference for \(\beta\) takes place in the quotient geometry
\[
    \LtwoT/\N_{K,\Lambda,T}
    \quad\text{or equivalently in}\quad
    \N_{K,\Lambda,T}^{\perp_T}.
\]
The projected moment is the inner product of two normal components:
\[
    \Psi_T(\beta,\Lambda)
    =\innerT{\bm M_{K,\Lambda}(\bm y-\beta\bm x)}
             {\bm M_{K,\Lambda}\bm w}.
\]

The empirical likelihood part adds a second geometry. For the scalar moment
\(g_t(\beta,\Lambda)\), define
\[
    C_T(\beta,\Lambda)
    =\left\{\pi\in\Delta_T:
      \sum_{t=1}^T\pi_tg_t(\beta,\Lambda)=0\right\},
\]
where \(\Delta_T=\{\pi_t\ge0,\sum_t\pi_t=1\}\). EL is the intersection of the
probability simplex with a moment hyperplane. The multiplier \(\lambda_T\) is
the normal coordinate to that hyperplane, and Wilks is the local quadratic
curvature of the logarithmic barrier around the uniform weight vector.

Finally, estimated partitions are compared in the topology the proof actually
uses, not in operator norm. For a set of relevant directions \(\mathcal D_T\),
define
\[
    d_{\mathcal D,T}(M_1,M_0)
    =\sup_{a,b\in\mathcal D_T}
      \left|T^{-1/2}a'(M_1-M_0)b\right|.
\]
The relevant directions are the few vectors that enter the score and derivative,
such as \(\bm\mu_0,\bm u,\bm x,\bm w\), and in the efficient theorem
\(\bm w^*\). The estimated-break proof requires
\[
    d_{\mathcal D,T}(\bm M_{K,\widehat\Lambda},\bm M_0)=o_p(1),
\]
not \(\|\bm M_{K,\widehat\Lambda}-\bm M_0\|_{op}=o_p(1)\).




\section{Paper-Facing Procedure}

The main theorem uses the null-imposed one-break proof object. For a tested
value \(\beta\), define
\[
    \bm y_\beta=\bm y-\beta\bm x.
\]
All partition comparisons below are made with this \(\bm y_\beta\). Thus, for a
\(p\)-value at \(H_0:\beta=\beta_0\), set \(\beta=\beta_0\). For a confidence
set, repeat the same null-imposed profiling at each candidate \(\beta\).

Let \(q\in\{0,1\}\). The value \(q=0\) is the stable-nuisance model and
\(q=1\) is the one-nuisance-break model. The score weight is
\[
    w_t=\omega(Z_t),
\]
where \(\omega\) is arbitrary within the admissible score-weight class:
boundedness or the required max-score condition, nondegenerate projected
variance, and the projected LLN/CLT assumptions are imposed on the resulting
residualized score. Bounded choices are convenient in simulations, but the
theorem does not require a particular functional form such as a hyperbolic
tangent.

\begin{quote}
\noindent\textbf{Algorithm 1: Null-imposed one-break profile-sieve EL.}
\begin{enumerate}[leftmargin=2em]
\item \textbf{Input.} Data \(\{Y_t,X_{t-1},W_{t-1},Z_t\}_{t=1}^T\), tested value
\(\beta\), sieve dimension \(K\), minimum segment length, candidate orders
\(q\in\{0,1\}\), penalty scale \(\kappa_T R_T\), and admissible score weight
\(w_t=\omega(Z_t)\).

\item \textbf{Search candidate nuisance subspaces.} For each \(q\in\{0,1\}\)
and each admissible partition \(\Lambda\in\mathcal L_q\), form the block-sieve
matrix \(\bm Q_{K,\Lambda}\), projection
\(\Pi_{K,\Lambda,T}\), and residual-maker
\(\bm M_{K,\Lambda}=I-\Pi_{K,\Lambda,T}\).

\item \textbf{Profile under the tested value.} Compute
\[
    \bm u_\Lambda(\beta)=\bm M_{K,\Lambda}(\bm y-\beta\bm x),
    \qquad
    \bm v_\Lambda=\bm M_{K,\Lambda}\bm w,
\]
and the null-imposed profile criterion
\[
    RSS_K(\beta,\Lambda)=\|\bm u_\Lambda(\beta)\|_2^2.
\]

\item \textbf{Select the partition for each order.} Set
\[
    \widehat\Lambda_q(\beta)
    \in\arg\min_{\Lambda\in\mathcal L_q}RSS_K(\beta,\Lambda).
\]

\item \textbf{Select the order.} Use the workbook penalty
\[
    IC(q;\beta)
    =RSS_K\{\beta,\widehat\Lambda_q(\beta)\}
     +q\,\kappa_T R_T,
\]
and set
\[
    \widehat q(\beta)\in\arg\min_{q\in\{0,1\}}IC(q;\beta),
    \qquad
    \widehat\Lambda(\beta)=\widehat\Lambda_{\widehat q(\beta)}(\beta).
\]

\item \textbf{Construct scalar EL scores.} With
\(\widehat M_\beta=\bm M_{K,\widehat\Lambda(\beta)}\), define
\[
    g_t(\beta,\widehat\Lambda)
    =\{\widehat M_\beta(\bm y-\beta\bm x)\}_t
     \{\widehat M_\beta\bm w\}_t.
\]

\item \textbf{Solve the scalar EL KKT equation.} Find \(\lambda(\beta)\) such
that
\[
    \sum_{t=1}^T
    \frac{g_t(\beta,\widehat\Lambda)}
         {1+\lambda(\beta)g_t(\beta,\widehat\Lambda)}
    =0,
    \qquad
    1+\lambda(\beta)g_t(\beta,\widehat\Lambda)>0.
\]

\item \textbf{Compute the EL ratio.}
\[
    \ell_T(\beta)
    =2\sum_{t=1}^T
      \log\{1+\lambda(\beta)g_t(\beta,\widehat\Lambda)\}.
\]
Reject \(H_0:\beta=\beta_0\) when \(\ell_T(\beta_0)\) exceeds the
\(\chi_1^2\) critical value.
\end{enumerate}
\end{quote}

\begin{remark}[What Algorithm 1 does not claim]
Algorithm 1 is the proof object used by the main theorem. It does not assert
that a particular computational shortcut, such as seeded binary segmentation or
narrowest-over-threshold search, is part of the statistical contribution.
Multiple-break searches can be used as exploratory diagnostics, but the
paper-facing theorem and implementation in this workbook are the
\(q\in\{0,1\}\) null-imposed procedure above.
\end{remark}

\begin{assumption}[Nuisance covariate, sieve, and break signal]
\label{ass:nuisance-covariate-sieve}
The following primitive conditions are sufficient for the nuisance-covariate
part of the high-level good event. They do not replace the good-event
assumptions below; they explain how the \(W\)-driven block sieve can satisfy
them.
\begin{enumerate}[leftmargin=2em]
\item \textbf{Predetermined design.} The lagged design variables are observed
before the innovation:
\[
    X_{t-1},W_{t-1}\in\mathcal F_{t-1},
    \qquad
    E(U_t\mid\mathcal F_{t-1})=0.
\]
The nuisance covariate is \(W_{t-1}\). The score weight \(w_t\) is a separate
object.

\item \textbf{Stable nuisance-covariate law.} The process \(W_t\) has a stable
marginal law \(P_W\) and satisfies the mixing or empirical-process conditions
needed for uniform laws over the searched sieve and break classes. A primitive
route is stationary alpha mixing with sufficiently fast decay.

\item \textbf{Sieve basis and envelope.} Let
\[
    P_K(W)=(p_1(W),\ldots,p_K(W))'.
\]
The basis includes an intercept and has envelope
\[
    \zeta_K=\sup_w \|P_K(w)\|.
\]
The growth of \(K\) and \(\zeta_K\) is slow enough for the projection and
local-break perturbation bounds.

\item \textbf{Uniform segment Gram stability.} For every admissible
partition \(\Lambda\in\mathcal L_T\), the block-sieve design satisfies
\[
    0<c\le
    \lambda_{\min}\{T^{-1}\bm Q_{K,\Lambda}'\bm Q_{K,\Lambda}\}
    \le
    \lambda_{\max}\{T^{-1}\bm Q_{K,\Lambda}'\bm Q_{K,\Lambda}\}
    \le C<\infty
\]
with probability tending to one.

\item \textbf{Regime-wise sieve approximation.} For each true regime
\(j=0,\ldots,q_0\), \(m_j\) belongs to a smoothness class approximable by
\(P_K\). There exists \(a_K\downarrow0\) such that
\[
    \max_{0\le j\le q_0}
    \inf_{c_j\in\mathbb R^K}
    \|m_j(W)-P_K(W)'c_j\|_{L^2(P_W)}
    \le a_K.
\]
The induced true-partition remainder
\(\bm r_{\mu,T}=\bm M_0\bm\mu_0\) is score-negligible:
\[
    \|\bm r_{\mu,T}\|_T=o_p(1),
    \qquad
    \sqrt T\,|\innerT{\bm r_{\mu,T}}{\bm M_0\bm w}|=o_p(1),
\]
\[
    P_T\{r_{\mu,T,t}^2(\bm M_0\bm w)_t^2\}=o_p(1).
\]

\item \textbf{Break spacing and break signal.} In the paper-facing one-break
case, \(q_0=1\) and the true break \(k_0\) satisfies
\[
    \min(k_0,T-k_0)\ge\epsilon T
\]
for some \(\epsilon>0\). The nuisance break strength is
\[
    \Delta_T^2
    =\|m_1(W)-m_0(W)\|_{L^2(P_W)}^2,
\]
and the RSS drift generated by this signal dominates searched stochastic
fluctuations outside the local window.

\item \textbf{Rate compatibility.} With
\[
    R_T=K+\log T+Ta_K^2,
    \qquad
    r_T=R_T/\Delta_T^2
\]
in the fixed-strength route, or the corresponding weak-break radius in the
detailed proof, assume
\[
    r_T/T\to0,
    \qquad
    \zeta_K\sqrt{Kr_T/T}\to0.
\]
For unknown order, the workbook penalty obeys
\[
    K\log T=o(\kappa_T R_T),
    \qquad
    \kappa_T R_T=o(T\Delta_T^2).
\]
\end{enumerate}
\end{assumption}

\section{Proof Architecture: Good-Event Geometry Before Probability}

The main-body proof is organized around the geometric dictionary just defined. Instead
of presenting many steps as stochastic lemmas first, it separates the argument
into two layers:
\[
\begin{array}{c}
    \text{deterministic geometry on a good event}\\[3pt]
    +\\[3pt]
    \text{stochastic proof that the good event}\\
    \text{has probability }1-o(1)
\end{array}
\]
The geometric layer is a finite-sample statement in \(L^2(P_T)\). It uses only
projection identities, KKT identities, distance comparisons, and deterministic
size inequalities. The stochastic layer proves that the deterministic hypotheses
of the geometric layer are satisfied with probability tending to one.

This is the main proof slogan:
\[
    \mathcal G_T
    \Longrightarrow
    \left\{
    \begin{array}{l}
    \text{moment bridge and consistency},\\
    \text{EL KKT quadratic reduction},\\
    \text{estimated-partition score and variance replacement},\\
    \text{Wilks reduction to a standardized score},\\
    \text{optional efficient-score linearization when }w=w^*.
    \end{array}
    \right.
\]
The supporting stochastic appendix proves \(\Pr(\mathcal G_T)\to1\), and the remaining stochastic
input is only the CLT/LLN for the oracle score.

\section{The Good Event}

Fix a deterministic sequence \(\delta_T\downarrow0\) and constants
\(0<c<C<\infty\). The good event is written as a set of geometric regularity
conditions:
\[
    \mathcal G_T
    =G_T^{\mathrm{iso}}
     \cap G_T^{\mathrm{approx}}
     \cap G_T^{\mathrm{angle}}
     \cap G_T^{\mathrm{local}}
     \cap G_T^{\mathrm{select}}
     \cap G_T^{\mathrm{convex}}.
\]
These names are not cosmetic. They identify which geometric fact each
stochastic lemma must certify.

\begin{enumerate}[leftmargin=2em]
\item \textbf{Approximate isometry: \(G_T^{\mathrm{iso}}\).} For all searched
partitions \(\Lambda\in\mathcal L_T\), the empirical Gram matrix of
\(\bm Q_{K,\Lambda}\) is nonsingular with eigenvalues in \([c,C]\). The empirical
inner products appearing in \(\Psi_T\), the derivative, and the variance differ
from their population analogues by at most \(\delta_T\), uniformly over the
needed finite-dimensional directions. This says the random \(L^2(P_T)\) geometry
is approximately isometric to the population geometry on the relevant subspaces.

\item \textbf{Small nuisance distance: \(G_T^{\mathrm{approx}}\).} With
\(\bm r_{\mu,T}=\bm M_0\bm\mu_0\),
\[
    \|\bm r_{\mu,T}\|_T\le \delta_T,
    \qquad
    \sqrt T\,|\innerT{\bm r_{\mu,T}}{\bm v_0}|\le \delta_T,
    \qquad
    P_T(r_{\mu,T,t}^2v_{0,t}^2)\le \delta_T.
\]
This says the true nuisance surface lies close to the true break-aware sieve
subspace, and its normal residue is negligible in the score direction.

\item \textbf{Angle and nondegeneracy: \(G_T^{\mathrm{angle}}\).} Let
\[
    A_T^0=T^{-1/2}\bm u'\bm v_0,
    \qquad
    B_T^0=T^{-1}\sum_{t=1}^T u_t^2v_{0,t}^2.
\]
The score is bounded at the local KKT scale, \(|A_T^0|\le C\), and the variance
is nondegenerate, \(B_T^0\ge c\). At the population level, the transversality
constant satisfies \(|\Gamma|\ge c\). This is the finite-sample version of the
statement that the residualized parametric direction has a nonzero angle with
the chosen residualized score direction.

\item \textbf{Local subspace stability: \(G_T^{\mathrm{local}}\).} The estimated
partition lies in the local window, \(\widehat\Lambda\in\mathcal N_T(r_T)\), and
throughout this window,
\[
    \sup_{\Lambda\in\mathcal N_T(r_T)}
    \sqrt T\,|\Psi_T(\beta_0,\Lambda)-\Psi_T(\beta_0,\Lambda_0)|\le\delta_T,
\]
\[
    \sup_{\Lambda\in\mathcal N_T(r_T)}
    |V_T(\Lambda)-V_T(\Lambda_0)|\le\delta_T,
\]
\[
    \sup_{\Lambda\in\mathcal N_T(r_T)}\sup_{\beta\in\mathcal B}
    |\Psi_T(\beta,\Lambda)-\Psi_T(\beta,\Lambda_0)|\le\delta_T.
\]
Equivalently, the residual-makers are close in the directional topology needed
by the proof, not necessarily in operator norm.

\item \textbf{Subspace selection: \(G_T^{\mathrm{select}}\).} Outside the local
window, the projection-distance drift dominates the noise fluctuation:
\[
    \inf_{\Lambda\in\mathcal L_T\setminus\mathcal N_T(r_T)}
    \{RSS_K(\beta_0,\Lambda)-RSS_K(\beta_0,\Lambda_0)\}>0.
\]
For unknown order, every underfitted order has larger penalized criterion than
\(q_0\), and every overfitted order also has larger penalized criterion than
\(q_0\). This is the nearest-subspace selection event.

\item \textbf{Convex feasibility: \(G_T^{\mathrm{convex}}\).} The max condition
holds,
\[
    \max_t|g_t(\beta_0,\widehat\Lambda)|/\sqrt T\le\delta_T,
\]
and the scalar convex hull of
\(\{g_t(\beta_0,\widehat\Lambda):1\le t\le T\}\) contains zero in its interior.
This says the simplex-hyperplane intersection defining EL is locally feasible
and the logarithmic barrier expansion is valid.
\end{enumerate}

The supporting stochastic appendix certifies
\[
    \Pr(\mathcal G_T)\to1
\]
for suitable \(\delta_T\downarrow0\), constants, and rates.
\section{Geometry Theorem I: Consistency Is Deterministic on \texorpdfstring{\(\mathcal G_T\)}{GT}}

\begin{theorem}[Good-event consistency]
\label{thm:v2-good-event-consistency}
Assume the population geometry from the detailed finite-sample geometry gives
\[
    \Psi(\beta,\Lambda_0)=-(\beta-\beta_0)\Gamma,
    \qquad \Gamma\ne0.
\]
On \(\mathcal G_T\), if \(\widehat\beta_T\in\mathcal B\) satisfies
\[
    |\Psi_T(\widehat\beta_T,\widehat\Lambda)|\le\delta_T,
\]
then
\[
    |\widehat\beta_T-\beta_0|\le \frac{2\delta_T}{|\Gamma|}
\]
for all large \(T\). Hence consistency follows once
\(\Pr(\mathcal G_T)\to1\).
\end{theorem}

\begin{proof}
On \(\mathcal G_T\), the uniform moment bridge gives
\[
    \sup_{\beta\in\mathcal B}
    |\Psi_T(\beta,\widehat\Lambda)-\Psi(\beta,\Lambda_0)|\le\delta_T.
\]
Therefore
\[
    |\Psi(\widehat\beta_T,\Lambda_0)|
    \le |\Psi(\widehat\beta_T,\Lambda_0)
          -\Psi_T(\widehat\beta_T,\widehat\Lambda)|
       +|\Psi_T(\widehat\beta_T,\widehat\Lambda)|
    \le 2\delta_T.
\]
The population geometry gives
\(|\Psi(\widehat\beta_T,\Lambda_0)|=|\widehat\beta_T-\beta_0||\Gamma|\), which
proves the claim.
\end{proof}

\section{Geometry Theorem II: EL Is Convex KKT Plus a Size Event}

\begin{theorem}[Good-event EL quadratic reduction]
\label{thm:v2-good-event-el}
On \(\mathcal G_T\), the scalar EL dual root at
\((\beta_0,\widehat\Lambda)\) exists and satisfies
\[
    \lambda_T
    =\frac{P_Tg_t(\beta_0,\widehat\Lambda)}
           {P_Tg_t(\beta_0,\widehat\Lambda)^2}
      +o_{\mathcal G}(T^{-1/2}),
\]
where \(o_{\mathcal G}(\cdot)\) denotes a deterministic remainder controlled on
\(\mathcal G_T\). Moreover,
\[
    \ell_T(\beta_0,
    \widehat\Lambda)
    =\frac{T\{\Psi_T(\beta_0,\widehat\Lambda)\}^2}
           {V_T(\widehat\Lambda)}+o_{\mathcal G}(1).
\]
\end{theorem}

\begin{proof}
The EL problem is the finite-dimensional convex simplex problem from the detailed finite-sample geometry.
On \(\mathcal G_T\), the scalar convex hull contains zero in its interior, so the
KKT multiplier \(\lambda_T\) exists and satisfies
\[
    P_T\left\{\frac{g_t}{1+\lambda_Tg_t}\right\}=0,
    \qquad 1+\lambda_Tg_t>0.
\]
The same event gives \(P_Tg_t=O(T^{-1/2})\),
\(P_Tg_t^2\ge c/2\), and \(\max_t|g_t|/\sqrt T\le\delta_T\). The deterministic
KKT expansion from Lemma~\ref{lem:part2-scalar-el-proof-claim} then gives the
multiplier expansion and
\[
    2\sum_t\log(1+\lambda_Tg_t)
    =\frac{T(P_Tg_t)^2}{P_Tg_t^2}+o_{\mathcal G}(1).
\]
Since \(P_Tg_t=\Psi_T(\beta_0,\widehat\Lambda)\) and
\(P_Tg_t^2=V_T(\widehat\Lambda)\), the displayed formula follows.
\end{proof}

\section{Geometry Theorem III: Directional Break Perturbations}

The important v2 point is that the estimated-partition argument does not need
operator-norm convergence of \(\bm M_{K,\widehat\Lambda}\) to \(\bm M_0\).
Projection matrices for two nearby break partitions may differ substantially as
operators. What matters is their action on the few directions entering the
score:
\[
    \bm y_{\beta_0}^\circ=\bm\mu_0+\bm u,
    \qquad \bm x,
    \qquad \bm w.
\]

\begin{proposition}[Good-event estimated-partition bridge]
\label{prop:v2-good-event-break-bridge}
On \(\mathcal G_T\),
\[
    \sqrt T\{\Psi_T(\beta_0,\widehat\Lambda)
      -\Psi_T(\beta_0,\Lambda_0)\}=o_{\mathcal G}(1),
\]
\[
    V_T(\widehat\Lambda)-V_T(\Lambda_0)=o_{\mathcal G}(1),
\]
and
\[
    \sup_{\beta\in\mathcal B}
    |\Psi_T(\beta,\widehat\Lambda)-\Psi_T(\beta,\Lambda_0)|=o_{\mathcal G}(1).
\]
\end{proposition}

\begin{proof}
Because \(\widehat\Lambda\in\mathcal N_T(r_T)\) on \(\mathcal G_T\), the three
displayed inequalities are immediate from the local directional replacement
component of \(\mathcal G_T\). The geometric reason this is plausible is local
support: when a candidate break is \(h\) observations away from a true break,
the regime labels differ only on a boundary block of size \(h\). The nuisance
part is a projection residual supported on that block plus the sieve remainder,
and the stochastic part is the directional quantity
\(\bm u'(\bm M_{K,\Lambda}-\bm M_0)\bm w\). The supporting stochastic appendix proves these quantities are small with high probability; this proposition only uses the resulting good-event inequalities.
\end{proof}

\section{Geometry Theorem IV: Wilks Is a Two-Step Consequence}

\begin{theorem}[Good-event Wilks reduction]
\label{thm:v2-good-event-wilks-reduction}
On \(\mathcal G_T\),
\[
    \ell_T(\beta_0,\widehat\Lambda)
    =\frac{(A_T^0)^2}{B_T^0}+o_{\mathcal G}(1),
\]
where
\[
    A_T^0=T^{-1/2}\bm u'\bm v_0,
    \qquad
    B_T^0=T^{-1}\sum_{t=1}^T u_t^2v_{0,t}^2.
\]
Consequently, if the stochastic shell proves
\[
    A_T^0\Rightarrow N(0,V),
    \qquad B_T^0\to_p V>0,
\]
then
\[
    \ell_T(\beta_0,\widehat\Lambda)\Rightarrow\chi_1^2.
\]
\end{theorem}

\begin{proof}
Theorem~\ref{thm:v2-good-event-el} gives
\[
    \ell_T(\beta_0,\widehat\Lambda)
    =\frac{T\{\Psi_T(\beta_0,\widehat\Lambda)\}^2}
           {V_T(\widehat\Lambda)}+o_{\mathcal G}(1).
\]
The estimated-partition bridge replaces \(\widehat\Lambda\) by \(\Lambda_0\).
At the true partition,
\[
    \sqrt T\Psi_T(\beta_0,\Lambda_0)
    =T^{-1/2}\bm u'\bm v_0
      +T^{-1/2}\bm r_{\mu,T}'\bm v_0
    =A_T^0+o_{\mathcal G}(1),
\]
and
\[
    V_T(\Lambda_0)=B_T^0+o_{\mathcal G}(1).
\]
Substitution gives the displayed good-event reduction. The final
\(\chi_1^2\) conclusion is now purely stochastic Slutsky: geometry has reduced
EL to the square of the standardized oracle score.
\end{proof}

\section{Geometry Theorem V: Break Search Is Subspace Selection}

\begin{proposition}[Good-event break localization and order selection]
\label{prop:v2-good-event-rss-selection}
On the RSS separation component of \(\mathcal G_T\), the known-order RSS
estimator satisfies
\[
    \widehat\Lambda\in\mathcal N_T(r_T).
\]
On the unknown-order penalty component of \(\mathcal G_T\),
\[
    \widehat q=q_0.
\]
\end{proposition}

\begin{proof}
The finite-sample RSS identity from the detailed finite-sample geometry gives
\[
    RSS_K(\beta_0,\Lambda)
    =T\operatorname{dist}_T^2(\bm y-\beta_0\bm x,
      \N_{K,\Lambda,T}).
\]
Thus break search is selection among candidate nuisance subspaces. On the RSS
separation component of \(\mathcal G_T\), every partition outside the local
window has larger RSS than \(\Lambda_0\), so an RSS minimizer must lie inside
\(\mathcal N_T(r_T)\). For order selection, the same deterministic comparison
applies to the penalized criterion: underfitted orders have deterministic
omitted-break distance loss larger than any penalty saving, and overfitted
orders have penalty increases larger than any spurious distance gain. Therefore
\(\widehat q=q_0\) on the event.
\end{proof}

\section{Optional Efficiency}

The optional efficiency result also has a geometric good-event form. Add the
efficient model and the efficient weight
\[
    w^*=X/\sigma^2(A),
    \qquad A=(J_0,W),
\]
as in the optional theorem stated in the supporting material. Add the efficiency event \(\mathcal G_T^{\rm eff}\):
\[
    \left\|\bm M_0\bm w^*-
    \left(\frac{\widetilde X_1}{\sigma^2(A_1)},\ldots,
          \frac{\widetilde X_T}{\sigma^2(A_T)}\right)'\right\|_T
    \le\delta_T,
\]
\[
    T^{-1/2}|\bm r_{\mu,T}'\bm M_0\bm w^*|\le\delta_T,
\]
\[
    \sqrt T|\Psi_T^*(\beta_0,\widehat\Lambda)
      -\Psi_T^*(\beta_0,\Lambda_0)|\le\delta_T,
\]
and
\[
    \left|\innerT{\bm M_{K,\widehat\Lambda}\bm x}
      {\bm M_{K,\widehat\Lambda}\bm w^*}
      -\innerT{\bm M_0\bm x}{\bm M_0\bm w^*}\right|\le\delta_T,
    \qquad
    \innerT{\bm M_0\bm x}{\bm M_0\bm w^*}\to I_{\mathrm{eff}}>0.
\]

\begin{theorem}[Good-event efficient linearization]
\label{thm:v2-good-event-efficiency}
On \(\mathcal G_T\cap\mathcal G_T^{\rm eff}\), if
\[
    \Psi_T^*(\widehat\beta_T,\widehat\Lambda)=o_{\mathcal G}(T^{-1/2}),
\]
then
\[
    \sqrt T(\widehat\beta_T-\beta_0)
    =I_{\mathrm{eff}}^{-1}
      T^{-1/2}\sum_{t=1}^T
      \frac{U_t\{X_t-E(X_t\mid A_t)\}}{\sigma^2(A_t)}
      +o_{\mathcal G}(1).
\]
\end{theorem}

\begin{proof}
For fixed \(\Lambda\), \(\Psi_T^*(\beta,\Lambda)\) is affine in \(\beta\), with
slope
\[
    -\innerT{\bm M_{K,\Lambda}\bm x}{\bm M_{K,\Lambda}\bm w^*}.
\]
Expanding the root at \(\beta_0\) and using the derivative bridge gives
\[
    \sqrt T(\widehat\beta_T-\beta_0)
    =\left\{\innerT{\bm M_0\bm x}{\bm M_0\bm w^*}\right\}^{-1}
      \sqrt T\Psi_T^*(\beta_0,\Lambda_0)+o_{\mathcal G}(1).
\]
The efficient-score approximation event gives
\[
    \sqrt T\Psi_T^*(\beta_0,\Lambda_0)
    =T^{-1/2}\sum_{t=1}^T
      \frac{U_t\{X_t-E(X_t\mid A_t)\}}{\sigma^2(A_t)}+o_{\mathcal G}(1),
\]
and the denominator converges to \(I_{\mathrm{eff}}\). This proves the result.
The limiting normal distribution and efficiency bound are then supplied by the
stochastic shell: the efficient-score CLT and the standard Hilbert-space
information inequality.
\end{proof}

\section{What Remains Stochastic}

The trial main-body proof is geometric, but it does not pretend probability
has disappeared. The genuinely stochastic tasks are exactly these:
\begin{enumerate}[leftmargin=2em]
\item uniform empirical Gram and inner-product laws, which put the sample
projection geometry close to population geometry;
\item sieve approximation rates for \(\mu_0\), \(E(w\mid A)\), and, in the
efficient theorem, \(E(w^*\mid A)\);
\item RSS concentration, which proves the distance gap for wrong partitions is
not overwhelmed by noise;
\item break-date localization and unknown-order penalty probability statements;
\item martingale or mixing CLT for \(T^{-1/2}\bm u'\bm v_0\), and the matching
variance LLN;
\item max-score and convex-hull probability bounds for the scalar EL KKT root;
\item in the optional efficiency theorem, the efficient-score CLT and any
plug-in variance-weight replacement.
\end{enumerate}

Thus the main-body proof is:
\[
\begin{array}{c}
    \text{finite-sample geometry}\\
    \Longrightarrow\\
    \text{good-event finite-sample reductions}\\
    \Longrightarrow\\
    \text{supporting stochastic appendix proves }\Pr(\mathcal G_T)\to1\\
    \Longrightarrow\\
    \text{CLT gives Wilks or efficiency.}
\end{array}
\]
The stochastic lemmas certify the good event rather than carrying the main logical burden.
\section{Minimal Final Statement}

The complete proof is governed by one geometric object: the subspace-valued map
\[
    \Lambda\longmapsto \N_{K,\Lambda,T}\subset L^2(P_T).
\]
For a trial \(\beta\), profiling is the orthogonal projection of
\(\bm y-\beta\bm x\) onto this nuisance subspace. The profiled residual and the
score direction both live in the normal space
\(\N_{K,\Lambda,T}^{\perp_T}\), so the projected moment is
\[
    \Psi_T(\beta,
    \Lambda)=\innerT{\bm M_{K,\Lambda}(\bm y-\beta\bm x)}
                    {\bm M_{K,\Lambda}\bm w}.
\]
Population identification is the corresponding normal-space statement in
\(L^2(P)\): the nuisance disappears under \(M_0^P\), and transversality
\[
    \Gamma=\innerP{M_0^PX}{M_0^Pw}\ne0
\]
forces the population moment line to cross zero only at \(\beta_0\).

Break search is nearest-subspace selection:
\[
    RSS_K(\beta,\Lambda)
    =T\operatorname{dist}_T^2(\bm y-\beta\bm x,
      \N_{K,\Lambda,T}).
\]
A wrong break chooses the wrong subspace. A stable no-break nuisance space can
leave a nonzero normal component \(\bm M_{\mathrm{st}}\bm\mu_0\); false
predictivity is the nonzero angle between this omitted nuisance residue and the
residualized score direction.

EL adds convex geometry. The empirical likelihood feasible set is
\[
    C_T(\beta,\Lambda)=
    \left\{\pi\in\Delta_T:\sum_t\pi_tg_t(\beta,\Lambda)=0\right\},
\]
the intersection of the probability simplex with a moment hyperplane. The EL
multiplier is the normal coordinate to that hyperplane, and the Wilks expansion
is the local quadratic curvature of the logarithmic barrier around uniform
weights.

The stochastic lemmas certify the good event
\[
    \mathcal G_T
    =G_T^{\mathrm{iso}}
     \cap G_T^{\mathrm{approx}}
     \cap G_T^{\mathrm{angle}}
     \cap G_T^{\mathrm{local}}
     \cap G_T^{\mathrm{select}}
     \cap G_T^{\mathrm{convex}},
    \qquad \Pr(\mathcal G_T)\to1.
\]
On this event the proof is deterministic geometry: consistency follows from the
uniform moment bridge plus transversality; EL reduces to the square of the
standardized projected score; estimated breaks can replace true breaks because
\(\widehat M_T\) is close to \(M_0\) in the directional topology actually used by
the score, not in global operator norm. The remaining probability is the oracle
score CLT and variance LLN, giving
\[
    \ell_T(\beta_0,\widehat\Lambda)\Rightarrow\chi_1^2.
\]

The optional efficiency theorem is separate. Under the conditional Gaussian
location efficiency model and efficient weight \(w^*=X/\sigma^2(A)\), the same
normal-space projection gives
\[
    M_0^Pw^*=\frac{X-E(X\mid A)}{\sigma^2(A)}
\]
and hence the projected score equals the efficient score. With the efficient
sieve approximation and estimated-partition derivative bridge,
\[
    \sqrt T(\widehat\beta_T-\beta_0)
    =I_{\mathrm{eff}}^{-1}T^{-1/2}\sum_{t=1}^T
      \frac{U_t\{X_t-E(X_t\mid A_t)\}}{\sigma^2(A_t)}+
      o_p(1).
\]
Thus efficiency is a conditional add-on to the Wilks proof, not part of the
minimal validity argument.

\appendix
\part{Detailed Supporting Material}

\section*{Appendix Roadmap}
The appendix retains the detailed machinery behind the front proof. The first block gives the finite-sample Hilbert and convex KKT identities. The second block records the stochastic bounds that certify \(\Pr(\mathcal G_T)\to1\). The third block keeps the older merge derivation and the optional efficiency argument in full detail.
\addcontentsline{toc}{section}{Appendix Roadmap}
\section{Finite-Sample Hilbert Space: Viewing the Entire Model in \texorpdfstring{\(L^2(P_T)\)}{L2PT}}

\subsection{The empirical space}

Let
\[
    \Omega_T=\{1,\ldots,T\},
    \qquad
    P_T=T^{-1}\sum_{t=1}^T\delta_t.
\]
Every sample vector \(\bm a=(a_1,\ldots,a_T)'\) is a function
\(a_T:\Omega_T\to\R\), with \(a_T(t)=a_t\). The empirical inner product is
\[
    \innerT{\bm a}{\bm b}
    =\int a_T b_T\,dP_T
    =T^{-1}\sum_{t=1}^Ta_tb_t
    =T^{-1}\bm a'\bm b.
\]
Thus all variables in the sample are elements of \(\LtwoT\):
\[
    \bm y,\bm x,\bm w,\bm u,\bm\mu_0\in\LtwoT.
\]

\subsection{The full finite-sample model}

The break-aware model is the vector identity in \(\LtwoT\)
\begin{equation}
    \bm y=\beta_0\bm x+\bm\mu_0+\bm u,
    \label{eq:l2pt-model}
\end{equation}
where
\[
    \mu_{0,t}=\sum_{j=0}^{q_0}m_j(W_{t-1})\ind\{k_{j,0}<t\le k_{j+1,0}\}.
\]
So the full model says: \(\bm y\) lies in the affine set
\[
    \beta_0\bm x+\bm\mu_0+\bm u.
\]
The inferential problem is to remove \(\bm\mu_0\), not by estimating it as a
main parameter, but by projecting away its nuisance subspace.

\subsection{Candidate nuisance subspaces}

For a candidate partition \(\Lambda=(k_1,\ldots,k_q)\), define the block sieve
matrix
\[
    \bm Q_{K,\Lambda}
    =[
      \bm D_0(\Lambda)\bm P,
      \ldots,
      \bm D_q(\Lambda)\bm P
    ].
\]
The sample nuisance subspace is
\begin{equation}
    \N_{K,\Lambda,T}
    =\operatorname{col}(\bm Q_{K,\Lambda})\subset\LtwoT.
    \label{eq:sample-nuisance-space}
\end{equation}
An element of this space has the form
\[
    g_T(t)=\sum_{j=0}^q\ind\{t\in I_j(\Lambda)\}\bm P_K(W_{t-1})'c_j.
\]
This is the finite-sample version of a regime-specific nonparametric nuisance
function.

Let \(j_\Lambda(t)\) denote the regime index induced by \(\Lambda\). For a
regime-wise nuisance tuple \(h=(h_0,\ldots,h_q)\), define its stacked sample
vector by
\[
    \bm h_{\Lambda,T}
    =\big(h_{j_\Lambda(1)}(W_0),\ldots,
          h_{j_\Lambda(T)}(W_{T-1})\big)'.
\]
This vector is the empirical version of the candidate nuisance direction.

\begin{lemma}[Uniform finite-sieve approximation in \texorpdfstring{\(L^2(P_T)\)}{L2PT}]
\label{lem:sieve-decomposition}
Fix \((K,\Lambda,T)\) and a nuisance class \(\mathcal H\). For each
\(h\in\mathcal H\), let \(\bm h_{\Lambda,T}\in\LtwoT\) be the stacked nuisance
vector defined above. Define the finite-sieve approximation radius
\begin{equation}
    \rho_{K,T}(\mathcal H,\Lambda)
    =\sup_{h\in\mathcal H}
      \inf_{\gamma\in\R^{(q+1)K}}
      \|\bm h_{\Lambda,T}-\bm Q_{K,\Lambda}\gamma\|_T.
    \label{eq:rho-sieve-radius}
\end{equation}
If \(\rho_{K,T}(\mathcal H,\Lambda)=\op(1)\), then the desired uniform
\(\LtwoT\) approximation holds:
\begin{equation}
    \sup_{h\in\mathcal H}
    \|\bm h_{\Lambda,T}-\bm Q_{K,\Lambda}\gamma_{K,\Lambda}(h)\|_T
    =\sup_{h\in\mathcal H}\|\bm r_{K,\Lambda,T}(h)\|_T
    =\rho_{K,T}(\mathcal H,\Lambda)
    =\op(1).
    \label{eq:uniform-sieve-approximation}
\end{equation}
Here, for each \(h\),
\[
    \gamma_{K,\Lambda}(h)
    \in\operatorname*{argmin}_{\gamma\in\R^{(q+1)K}}
    \|\bm h_{\Lambda,T}-\bm Q_{K,\Lambda}\gamma\|_T^2,
\]
\[
    \bm h_{K,\Lambda,T}
    =\bm Q_{K,\Lambda}\gamma_{K,\Lambda}(h),
    \qquad
    \bm r_{K,\Lambda,T}(h)
    =\bm h_{\Lambda,T}-\bm h_{K,\Lambda,T}.
\]
The ambient space is \(\LtwoT\cong\R^T\), so \(\bm h_{\Lambda,T}\),
\(\bm h_{K,\Lambda,T}\), and \(\bm r_{K,\Lambda,T}(h)\) are all \(T\times1\)
vectors. The design matrix and coefficient vector have dimensions
\[
    \bm Q_{K,\Lambda}\in\R^{T\times(q+1)K},
    \qquad
    \gamma_{K,\Lambda}(h)\in\R^{(q+1)K}.
\]
The finite nuisance space has dimension
\[
    d_{K,\Lambda,T}=\operatorname{rank}(\bm Q_{K,\Lambda})\le(q+1)K,
\]
and
\begin{equation}
    \bm h_{\Lambda,T}
    =\bm Q_{K,\Lambda}\gamma_{K,\Lambda}(h)
     +\bm r_{K,\Lambda,T}(h),
    \label{eq:sieve-decomposition}
\end{equation}
where
\[
    \bm Q_{K,\Lambda}\gamma_{K,\Lambda}(h)\in\N_{K,\Lambda,T}\subset\LtwoT,
    \qquad
    \bm r_{K,\Lambda,T}(h)\in\N_{K,\Lambda,T}^{\perp_T}\subset\LtwoT.
\]
Thus \(\bm r_{K,\Lambda,T}(h)\) is a \(T\times1\) approximation-residual vector
living in a \((T-d_{K,\Lambda,T})\)-dimensional orthogonal complement, and
\[
    \bm Q_{K,\Lambda}'\bm r_{K,\Lambda,T}(h)=0.
\]
\end{lemma}

\begin{proof}
The column space \(\N_{K,\Lambda,T}=\operatorname{col}(\bm Q_{K,\Lambda})\) is a
finite-dimensional closed subspace of \(\LtwoT\). Therefore the Hilbert
projection theorem gives the closest finite-sieve approximation to every
\(\bm h_{\Lambda,T}\):
\[
    \Pi_{K,\Lambda,T}\bm h_{\Lambda,T}
    \in\N_{K,\Lambda,T}.
\]
Because every element of \(\N_{K,\Lambda,T}\) is \(\bm Q_{K,\Lambda}\gamma\), the
projection can be written as
\[
    \Pi_{K,\Lambda,T}\bm h_{\Lambda,T}
    =\bm Q_{K,\Lambda}\gamma_{K,\Lambda}(h)
\]
for any least-squares minimizer \(\gamma_{K,\Lambda}(h)\). If
\(\bm Q_{K,\Lambda}\) is rank deficient, the coefficient vector need not be
unique, but the fitted vector is unique.

Define
\[
    \bm r_{K,\Lambda,T}(h)
    =\bm h_{\Lambda,T}-\Pi_{K,\Lambda,T}\bm h_{\Lambda,T}.
\]
The projection theorem gives
\[
    \innerT{\bm r_{K,\Lambda,T}(h)}{\bm v}=0
    \qquad\text{for every }\bm v\in\N_{K,\Lambda,T},
\]
which implies \(\bm Q_{K,\Lambda}'\bm r_{K,\Lambda,T}(h)=0\). This proves the
orthogonal decomposition \eqref{eq:sieve-decomposition} and the location
\(\bm r_{K,\Lambda,T}(h)\in\N_{K,\Lambda,T}^{\perp_T}\).

By the definition of orthogonal projection,
\[
    \|\bm r_{K,\Lambda,T}(h)\|_T
    =\inf_{\gamma\in\R^{(q+1)K}}
      \|\bm h_{\Lambda,T}-\bm Q_{K,\Lambda}\gamma\|_T.
\]
Taking the supremum over \(h\in\mathcal H\) gives exactly
\eqref{eq:uniform-sieve-approximation}. Hence the mathematical task behind this
lemma is to verify \(\rho_{K,T}(\mathcal H,\Lambda)=\op(1)\) from the chosen
basis, smoothness class, and growth rate of \(K\). An arbitrary nuisance
direction is not assumed to be exactly spanned by the sieve basis; it is
approximated by its finite-sieve projection, and the leftover is
\(\bm r_{K,\Lambda,T}(h)\).
\end{proof}
\subsection{Orthogonal projection}

Let \(\Pi_{K,\Lambda,T}\) be the \(\LtwoT\)-orthogonal projection onto
\(\N_{K,\Lambda,T}\). In coordinates,
\[
    \Pi_{K,\Lambda,T}
    =\bm Q_{K,\Lambda}(\bm Q_{K,\Lambda}'\bm Q_{K,\Lambda})^\dagger
      \bm Q_{K,\Lambda}'.
\]
The residual-maker is
\begin{equation}
    \bm M_{K,\Lambda}=I-\Pi_{K,\Lambda,T}.
    \label{eq:sample-residual-maker}
\end{equation}
It satisfies the projection identities
\[
    \bm M_{K,\Lambda}'=\bm M_{K,\Lambda},
    \qquad
    \bm M_{K,\Lambda}^2=\bm M_{K,\Lambda},
    \qquad
    \bm Q_{K,\Lambda}'\bm M_{K,\Lambda}=0.
\]

\begin{lemma}[Projection identities for the residual-maker]
\label{lem:projection-identities}
Let \(Q=\bm Q_{K,\Lambda}\), \(G=Q'Q\),
\[
    \Pi=QG^\dagger Q',
    \qquad
    M=I-\Pi,
\]
where \(G^\dagger\) is the Moore-Penrose inverse. Then \(\Pi\) and \(M\) are
symmetric and idempotent, and \(M\) annihilates the sieve nuisance space:
\[
    M'=M,
    \qquad
    M^2=M,
    \qquad
    MQ=0,
    \qquad
    Q'M=0.
\]
\end{lemma}

\begin{proof}
Since \(G=Q'Q\) is symmetric, its Moore-Penrose inverse \(G^\dagger\) is also
symmetric. Hence
\[
    \Pi'=\{QG^\dagger Q'\}'=QG^\dagger Q'=\Pi.
\]
For idempotence, use the Moore-Penrose identity
\(G^\dagger GG^\dagger=G^\dagger\):
\[
    \Pi^2
    =QG^\dagger Q'QG^\dagger Q'
    =QG^\dagger GG^\dagger Q'
    =QG^\dagger Q'
    =\Pi.
\]
Therefore
\[
    M'=I-\Pi'=M,
    \qquad
    M^2=(I-\Pi)^2=I-2\Pi+\Pi^2=I-\Pi=M.
\]
It remains to show that \(M\) kills the columns of \(Q\). Since
\(\mathcal R(G)=\mathcal R(Q')\) and \(G^\dagger G\) is the orthogonal projector
onto \(\mathcal R(G)\),
\[
    QG^\dagger G=Q.
\]
Thus
\[
    \Pi Q=QG^\dagger Q'Q=QG^\dagger G=Q,
    \qquad
    MQ=(I-\Pi)Q=0.
\]
Finally, because \(M\) is symmetric, \(Q'M=(MQ)'=0\). This proves all projection
identities.
\end{proof}

\section{Profiling Is Orthogonal Projection}

For a trial value \(\beta\), define the partial outcome
\[
    \bm y_\beta^\circ=\bm y-\beta\bm x.
\]
It is the original outcome \(\bm y\) after subtracting the trial linear component
\(\beta\bm x\).

Profiling the nuisance means finding the closest nuisance vector to the partial
outcome \(\bm y_\beta^\circ\) in \(\N_{K,\Lambda,T}\):
\[
    \widehat{\bm\mu}_{\beta,\Lambda}
    =\arg\min_{\bm g\in\N_{K,\Lambda,T}}\|\bm y_\beta^\circ-\bm g\|_T^2.
\]
By Hilbert projection,
\[
    \widehat{\bm\mu}_{\beta,\Lambda}=\Pi_{K,\Lambda,T}\bm y_\beta^\circ,
    \qquad
    \widehat{\bm u}(\beta,\Lambda)=\bm y_\beta^\circ-\widehat{\bm\mu}_{\beta,\Lambda}
    =\bm M_{K,\Lambda}\bm y_\beta^\circ.
\]

\begin{lemma}[Profile normal equation]
For every \(\bm h\in\N_{K,\Lambda,T}\),
\[
    \innerT{\widehat{\bm u}(\beta,\Lambda)}{\bm h}=0.
\]
Equivalently,
\[
    \bm Q_{K,\Lambda}'\widehat{\bm u}(\beta,\Lambda)=0.
\]
\end{lemma}

\begin{proof}
This is the KKT condition for a least-squares projection in \(\LtwoT\). If
\(\bm h=\bm Q_{K,\Lambda}\bm a\), then
\[
    \innerT{\widehat{\bm u}}{\bm h}
    =T^{-1}\widehat{\bm u}'\bm Q_{K,\Lambda}\bm a
    =T^{-1}(\bm Q_{K,\Lambda}'\widehat{\bm u})'\bm a=0.
\]
\end{proof}

\subsection{Break search as subspace selection}

For fixed \(\beta\), the partial outcome is
\[
    \bm y_\beta^\circ=\bm y-\beta\bm x.
\]
For a candidate partition \(\Lambda=(k_1,\ldots,k_q)\), set \(k_0=0\),
\(k_{q+1}=T\), and
\[
    I_j(\Lambda)=\{t:k_j<t\le k_{j+1}\}.
\]
The block sieve matrix is
\[
    \bm Q_{K,\Lambda}
    =[\bm D_0(\Lambda)\bm P,\ldots,\bm D_q(\Lambda)\bm P].
\]
For \(c=(c_0',\ldots,c_q')'\), the fitted nuisance vector is
\(\bm Q_{K,\Lambda}c\), with component
\[
    (\bm Q_{K,\Lambda}c)_t
    =\bm P_K(W_{t-1})'c_j
    \qquad\text{if }t\in I_j(\Lambda).
\]
Thus minimizing over \(c\) is exactly the same as choosing a separate sieve
coefficient vector \(c_j\) on each segment.

\begin{lemma}[Profile RSS is a Hilbert distance]
\label{lem:rss-hilbert-distance}
For fixed \((\beta,\Lambda)\), define the segment cost
\[
    C_K(s,e;\beta)
    =\min_{c\in\R^K}
      \sum_{t=s}^e
      \{y_t-\beta x_t-\bm P_K(W_{t-1})'c\}^2,
\]
and the profile RSS
\[
    RSS_K(\beta,\Lambda)
    =\sum_{j=0}^q C_K(k_j+1,k_{j+1};\beta).
\]
Then
\begin{equation}
\begin{aligned}
    RSS_K(\beta,\Lambda)
    &=\min_{c\in\R^{(q+1)K}}
      \|\bm y_\beta^\circ-\bm Q_{K,\Lambda}c\|_2^2 \\
    &=\inf_{\bm g\in\N_{K,\Lambda,T}}
      \|\bm y_\beta^\circ-\bm g\|_2^2 \\
    &=T\operatorname{dist}_T^2
      (\bm y_\beta^\circ,\N_{K,\Lambda,T}) \\
    &=T\|\bm M_{K,\Lambda}\bm y_\beta^\circ\|_T^2.
\end{aligned}
\label{eq:rss-distance-equivalence}
\end{equation}
\end{lemma}

\begin{proof}
Because each segment \(I_j(\Lambda)\) has its own coefficient vector \(c_j\),
\[
\begin{aligned}
    RSS_K(\beta,\Lambda)
    &=\min_{c_0,\ldots,c_q}
      \sum_{j=0}^q\sum_{t\in I_j(\Lambda)}
      \{y_t-\beta x_t-\bm P_K(W_{t-1})'c_j\}^2 \\
    &=\min_{c\in\R^{(q+1)K}}
      \|\bm y-\beta\bm x-\bm Q_{K,\Lambda}c\|_2^2 \\
    &=\min_{c\in\R^{(q+1)K}}
      \|\bm y_\beta^\circ-\bm Q_{K,\Lambda}c\|_2^2.
\end{aligned}
\]
Since \(\N_{K,\Lambda,T}=\operatorname{col}(\bm Q_{K,\Lambda})\), every
\(\bm g\in\N_{K,\Lambda,T}\) can be written as
\(\bm g=\bm Q_{K,\Lambda}c\). Therefore
\[
    \min_c\|\bm y_\beta^\circ-\bm Q_{K,\Lambda}c\|_2^2
    =\inf_{\bm g\in\N_{K,\Lambda,T}}
      \|\bm y_\beta^\circ-\bm g\|_2^2.
\]
The Euclidean norm and the empirical Hilbert norm satisfy
\(\|\bm a\|_2^2=T\|\bm a\|_T^2\). Hence
\[
    \inf_{\bm g\in\N_{K,\Lambda,T}}
      \|\bm y_\beta^\circ-\bm g\|_2^2
    =T\operatorname{dist}_T^2(\bm y_\beta^\circ,\N_{K,\Lambda,T}).
\]
Finally, \(\bm M_{K,\Lambda}\) is the residual-maker for the orthogonal
projection onto \(\N_{K,\Lambda,T}\), so
\[
    \operatorname{dist}_T^2(\bm y_\beta^\circ,\N_{K,\Lambda,T})
    =\|\bm M_{K,\Lambda}\bm y_\beta^\circ\|_T^2.
\]
This proves \eqref{eq:rss-distance-equivalence}.
\end{proof}

Consequently, for known order \(q\), the break estimator is
\[
\begin{aligned}
    \widehat\Lambda_q(\beta)
    &=\arg\min_{\Lambda\in\mathcal L_q(\epsilon)} RSS_K(\beta,\Lambda) \\
    &=\arg\min_{\Lambda\in\mathcal L_q(\epsilon)}
      \operatorname{dist}_T^2(\bm y_\beta^\circ,\N_{K,\Lambda,T}).
\end{aligned}
\]
Unknown order adds a penalty to the same distance:
\[
    \widehat q(\beta)
    =\arg\min_{0\le q\le\bar q}
     \left\{
       T\operatorname{dist}_T^2
       (\bm y_\beta^\circ,\N_{K,\widehat\Lambda_q(\beta),T})
       +\operatorname{pen}_T(q,K)
     \right\}.
\]
Thus break search is subspace selection: each candidate \(\Lambda\) creates a
candidate nuisance subspace \(\N_{K,\Lambda,T}\), and RSS picks the subspace
closest to the partial outcome \(\bm y_\beta^\circ\).
\section{Residualized Score Direction}

The raw score weight \(\bm w\in\LtwoT\) may contain nuisance directions. The
geometric score direction is its nuisance-orthogonal component
\[
    \bm v_\Lambda=\bm w^c(\Lambda)=\bm M_{K,\Lambda}\bm w.
\]
Then
\[
    \bm v_\Lambda\in\N_{K,\Lambda,T}^{\perp_T},
    \qquad
    \innerT{\bm v_\Lambda}{\bm h}=0
    \quad \forall \bm h\in\N_{K,\Lambda,T}.
\]

The sample profiled moment is the Hilbert inner product
\begin{equation}
    \Psi_T(\beta,\Lambda)
    =\innerT{\bm M_{K,\Lambda}(\bm y-\beta\bm x)}
      {\bm M_{K,\Lambda}\bm w}.
    \label{eq:sample-profile-moment}
\end{equation}
The scaled score is
\[
    S_T(\beta,\Lambda)=\sqrt T\,\Psi_T(\beta,\Lambda).
\]

\begin{lemma}[Finite-sample nuisance cancellation]
If \(\bm g\in\N_{K,\Lambda,T}\), then
\[
    \innerT{\bm M_{K,\Lambda}\bm g}{\bm M_{K,\Lambda}\bm w}=0.
\]
\end{lemma}

\begin{proof}
Because \(\bm g\in\N_{K,\Lambda,T}\), \(\bm M_{K,\Lambda}\bm g=0\). Hence the
inner product is zero. Equivalently,
\((\bm Qc)'\bm M_{K,\Lambda}\bm w=0\) for every \(c\).
\end{proof}

\section{Known True Partition: Wilks Logic in Geometry}

Set \(\Lambda=\Lambda_0\), and abbreviate
\[
    \bm Q_0=\bm Q_{K,\Lambda_0},
    \qquad
    \bm M_0=\bm M_{K,\Lambda_0},
    \qquad
    \bm v_0=\bm M_0\bm w.
\]
Decompose the true nuisance into its finite-sieve projection and approximation
residual:
\[
    \bm\mu_0=\bm Q_0\bm c_0+\bm r_0,
    \qquad
    \bm Q_0\bm c_0\in\N_{K,\Lambda_0,T},
    \qquad
    \bm r_0\in\N_{K,\Lambda_0,T}^{\perp_T}.
\]

\begin{lemma}[True-partition score decomposition]
\label{lem:true-partition-score-decomposition}
At \(\beta=\beta_0\) and \(\Lambda=\Lambda_0\), the projected score satisfies
\begin{equation}
    S_T(\beta_0,\Lambda_0)
    =T^{-1/2}\bm u'\bm v_0
     +T^{-1/2}\bm r_0'\bm v_0.
    \label{eq:true-partition-score-decomposition}
\end{equation}
The first term is the stochastic score. The second term is the sieve
approximation-bias term, and this is the term controlled by approximation
rates:
\begin{equation}
    T^{-1/2}\bm r_0'\bm v_0
    =T^{-1/2}\bm r_0'\bm M_0\bm w
    =\op(1).
    \label{eq:true-partition-approximation-bias}
\end{equation}
\end{lemma}

\begin{proof}
At the true parameter, the partial outcome is
\[
    \bm y_{\beta_0}^\circ
    =\bm y-\beta_0\bm x
    =\bm\mu_0+\bm u
    =\bm Q_0\bm c_0+\bm r_0+\bm u.
\]
Apply the true-partition residual-maker \(\bm M_0\). Since \(\bm Q_0\bm c_0\) is
in \(\N_{K,\Lambda_0,T}\), Lemma~\ref{lem:projection-identities} gives
\(\bm M_0\bm Q_0\bm c_0=0\). Therefore
\[
    \bm M_0\bm y_{\beta_0}^\circ
    =\bm M_0\bm u+\bm M_0\bm r_0.
\]
The projected score is
\[
\begin{aligned}
    S_T(\beta_0,\Lambda_0)
    &=\sqrt T\innerT{\bm M_0\bm y_{\beta_0}^\circ}{\bm M_0\bm w} \\
    &=T^{-1/2}(\bm M_0\bm u+\bm M_0\bm r_0)'\bm v_0.
\end{aligned}
\]
Because \(\bm M_0\) is symmetric and idempotent and \(\bm v_0=\bm M_0\bm w\),
we have \(\bm M_0\bm v_0=\bm v_0\). Hence
\[
    (\bm M_0\bm u)'\bm v_0=\bm u'\bm v_0,
    \qquad
    (\bm M_0\bm r_0)'\bm v_0=\bm r_0'\bm v_0.
\]
Substituting these two identities gives
\eqref{eq:true-partition-score-decomposition}. The term
\(T^{-1/2}\bm u'\bm v_0\) is handled by the score CLT and variance consistency.
The term \(T^{-1/2}\bm r_0'\bm v_0\) is the only approximation-rate term in this
decomposition, because it is the only term containing the sieve residual
\(\bm r_0\).
\end{proof}
\section{Omitted Breaks: False Predictivity as a Projection Angle}

Let \(\N_{K,0,T}\) be the stable/no-break nuisance space and let
\(\bm M_{\mathrm{st}}\) be the residual-maker for orthogonal projection onto
\(\N_{K,0,T}\). Define
\[
    \bm v_{\mathrm{st}}=\bm M_{\mathrm{st}}\bm w,
    \qquad
    S_T^{\mathrm{st}}(\beta)
    =\sqrt T\innerT{\bm M_{\mathrm{st}}(\bm y-\beta\bm x)}
                         {\bm M_{\mathrm{st}}\bm w}.
\]

\begin{proposition}[Omitted-break projection drift]
\label{prop:omitted-break-projection-drift}
At the true parameter \(\beta_0\),
\begin{equation}
    S_T^{\mathrm{st}}(\beta_0)
    =
    T^{-1/2}\bm u'\bm v_{\mathrm{st}}
    +B_T,
    \qquad
    B_T=T^{-1/2}\bm\mu_0'\bm v_{\mathrm{st}}.
    \label{eq:stable-score-drift-decomposition}
\end{equation}
Equivalently,
\begin{equation}
    B_T
    =
    \sqrt T
    \innerT{\bm M_{\mathrm{st}}\bm\mu_0}
           {\bm M_{\mathrm{st}}\bm w}.
    \label{eq:stable-score-drift-inner-product}
\end{equation}
If
\[
    \innerT{\bm M_{\mathrm{st}}\bm\mu_0}
           {\bm M_{\mathrm{st}}\bm w}
    \to b\ne0,
\]
then \(B_T/\sqrt T\to b\), so the stable score has deterministic drift of
order \(\sqrt T\) at the true parameter.
\end{proposition}

\begin{proof}
At \(\beta_0\),
\[
    \bm y-\beta_0\bm x=\bm\mu_0+\bm u.
\]
Applying the stable residual-maker gives
\[
    \bm M_{\mathrm{st}}(\bm y-\beta_0\bm x)
    =
    \bm M_{\mathrm{st}}\bm u
    +
    \bm M_{\mathrm{st}}\bm\mu_0.
\]
Therefore
\[
\begin{aligned}
    S_T^{\mathrm{st}}(\beta_0)
    &=
    T^{-1/2}
    (\bm M_{\mathrm{st}}\bm u+\bm M_{\mathrm{st}}\bm\mu_0)'
    \bm v_{\mathrm{st}}  \\
    &=
    T^{-1/2}\bm u'\bm v_{\mathrm{st}}
    +
    T^{-1/2}\bm\mu_0'\bm v_{\mathrm{st}},
\end{aligned}
\]
because \(\bm M_{\mathrm{st}}\) is symmetric and idempotent and
\(\bm v_{\mathrm{st}}=\bm M_{\mathrm{st}}\bm w\). This proves
\eqref{eq:stable-score-drift-decomposition}. The same symmetry and idempotence
give
\[
    T^{-1/2}\bm\mu_0'\bm v_{\mathrm{st}}
    =
    T^{-1/2}\bm\mu_0'\bm M_{\mathrm{st}}\bm w
    =
    \sqrt T
    \innerT{\bm M_{\mathrm{st}}\bm\mu_0}
           {\bm M_{\mathrm{st}}\bm w},
\]
which proves \eqref{eq:stable-score-drift-inner-product}. If the empirical
inner product converges to \(b\ne0\), the drift is \(B_T=\sqrt T\,b+o(\sqrt T)\).
\end{proof}

The same statement can be read as an angle condition. On the stable normal
space, define the empirical cosine
\[
    \cos_T(\bm M_{\mathrm{st}}\bm\mu_0,\bm v_{\mathrm{st}})
    =\frac{\innerT{\bm M_{\mathrm{st}}\bm\mu_0}{\bm v_{\mathrm{st}}}}
           {\|\bm M_{\mathrm{st}}\bm\mu_0\|_T\|\bm v_{\mathrm{st}}\|_T},
\]
when both norms are nonzero. Then
\[
    B_T=\sqrt T\,
    \|\bm M_{\mathrm{st}}\bm\mu_0\|_T\,\|\bm v_{\mathrm{st}}\|_T\,
    \cos_T(\bm M_{\mathrm{st}}\bm\mu_0,\bm v_{\mathrm{st}}).
\]
Thus false predictivity is a nonzero angle between the unremoved nuisance
residue and the residualized score direction. The correct break-aware projection
is needed because it makes the nuisance residue \(\bm M_0\bm\mu_0\) small in the
relevant normal space.
As a toy example, suppose there are no covariates and the stable nuisance space
contains only constants:
\[
    \N_{K,0,T}=\operatorname{span}\{\bm 1_T\}.
\]
Let the true nuisance contain one level shift,
\[
    \mu_{0,t}=\delta\ind\{t>k_0\},
    \qquad
    \pi_T=T^{-1}\sum_{t=1}^T\ind\{t>k_0\}.
\]
The stable projection removes only the global mean. Hence
\[
    (\bm M_{\mathrm{st}}\bm\mu_0)_t
    =
    \delta\{\ind(t>k_0)-\pi_T\}.
\]
If the residualized score direction is aligned with this step, for instance
\(\bm M_{\mathrm{st}}\bm w=\bm M_{\mathrm{st}}\bm\mu_0\), then
\[
    \innerT{\bm M_{\mathrm{st}}\bm\mu_0}
           {\bm M_{\mathrm{st}}\bm w}
    =
    \|\bm M_{\mathrm{st}}\bm\mu_0\|_T^2
    =
    \delta^2\pi_T(1-\pi_T).
\]
When \(\delta\ne0\) and \(\pi_T\to\pi\in(0,1)\), the limit is
\(\delta^2\pi(1-\pi)\ne0\). The stable score therefore reports a false signal
even at \(\beta_0\).

The correct break-aware block sieve changes the projection target. Under the
true partition,
\[
    \bm\mu_0=\bm Q_0\bm c_0+\bm r_0,
    \qquad
    \bm Q_0\bm c_0\in\N_{K,\Lambda_0,T},
    \qquad
    \|\bm r_0\|_T=o_p(1),
\]
so \(\bm M_0\bm\mu_0=\bm r_0\) is small. Under the stable space,
\(\bm M_{\mathrm{st}}\bm\mu_0\) need not be small. The block-sieve projection
removes the low-frequency break component before testing \(\beta\); otherwise
omitted nuisance breaks can enter the score as false predictivity.
\section{Estimated Partition}

The estimated break procedure chooses a random subspace
\[
    \widehat\N_T=\N_{K,\widehat\Lambda,T},
    \qquad
    \widehat M_T=\bm M_{K,\widehat\Lambda}.
\]
The feasible score is
\[
    \widehat S_T(\beta_0)
    =\sqrt T\innerT{\widehat M_T(\bm y-\beta_0\bm x)}{\widehat M_T\bm w}.
\]
The true-partition score is
\[
    S_T^0(\beta_0)
    =\sqrt T\innerT{\bm M_0(\bm y-\beta_0\bm x)}{\bm M_0\bm w},
    \qquad
    \bm M_0=\bm M_{K,\Lambda_0}.
\]
Because each residual-maker is symmetric and idempotent,
\[
    \widehat S_T(\beta_0)=T^{-1/2}(\bm\mu_0+\bm u)'\widehat M_T\bm w,
    \qquad
    S_T^0(\beta_0)=T^{-1/2}(\bm\mu_0+\bm u)'\bm M_0\bm w.
\]
Therefore the exact estimated-partition score error is
\begin{equation}
\begin{aligned}
    \widehat S_T(\beta_0)-S_T^0(\beta_0)
    &=T^{-1/2}\bm\mu_0'(\widehat M_T-\bm M_0)\bm w \\
    &\quad+T^{-1/2}\bm u'(\widehat M_T-\bm M_0)\bm w.
\end{aligned}
\label{eq:estimated-partition-score-decomposition}
\end{equation}
The first term is local sieve-bias perturbation. The second term is local noise
perturbation. Thus the estimated-partition Wilks argument needs
\begin{equation}
    T^{-1/2}\bm\mu_0'(\widehat M_T-\bm M_0)\bm w=\op(1),
    \qquad
    T^{-1/2}\bm u'(\widehat M_T-\bm M_0)\bm w=\op(1).
    \label{eq:estimated-partition-score-pieces}
\end{equation}
For the variance, the corresponding requirement is
\begin{equation}
    \widehat V_T-V_T^0=\op(1),
    \qquad
    \widehat V_T=P_Tg_t(\beta_0,\widehat\Lambda)^2,
    \qquad
    V_T^0=P_Tg_t(\beta_0,\Lambda_0)^2.
    \label{eq:estimated-partition-variance-replacement}
\end{equation}
Equations \eqref{eq:estimated-partition-score-decomposition}--\eqref{eq:estimated-partition-variance-replacement}
are the exact bridge from estimated breaks to the true-partition geometry.

The reason this is weaker than operator-norm convergence is visible from
\eqref{eq:estimated-partition-score-decomposition}. We do not need
\(\|\widehat M_T-\bm M_0\|_{op}\to0\). We only need the difference to be small
when applied to the score direction \(\bm w\) and paired with \(\bm\mu_0\) and
\(\bm u\), plus the analogous second-moment replacement. In the directional
topology of the geometric dictionary, this is a statement that
\[
    d_{\mathcal D,T}(\widehat M_T,\bm M_0)=o_p(1)
\]
for a direction set \(\mathcal D_T\) containing the vectors that appear in the
score, variance, and derivative. This is the topology induced by the theorem,
not an after-the-fact weakening of an operator-norm claim.

The localization proof uses the RSS distance identity from
Lemma~\ref{lem:rss-hilbert-distance}. At \(\beta_0\), set
\[
    \Delta_\Lambda=\bm M_{K,\Lambda}-\bm M_0,
    \qquad
    \bm y_{\beta_0}^\circ=\bm\mu_0+\bm u.
\]
Then
\begin{equation}
\begin{aligned}
    RSS_K(\beta_0,\Lambda)-RSS_K(\beta_0,\Lambda_0)
    &=(\bm y_{\beta_0}^\circ)'\Delta_\Lambda\bm y_{\beta_0}^\circ \\
    &=\bm\mu_0'\Delta_\Lambda\bm\mu_0
      +2\bm\mu_0'\Delta_\Lambda\bm u
      +\bm u'\Delta_\Lambda\bm u.
\end{aligned}
\label{eq:rss-localization-decomposition}
\end{equation}
The deterministic term is the geometric drift from choosing the wrong nuisance
subspace. The two stochastic terms must be uniformly smaller than that drift on
wrong partitions. Part II states the primitive concentration conditions and then
uses \eqref{eq:rss-localization-decomposition} to obtain break localization and
order selection.
\section{EL as a Finite-Sample Convex KKT Problem}

For fixed \((K,\Lambda,\beta)\), everything in this section is finite-sample
algebra in \(\LtwoT\). Define the profiled residual, score direction, and scalar
moment array
\[
    \bm r_\beta=\bm M_{K,\Lambda}(\bm y-\beta\bm x),
    \qquad
    \bm a=\bm M_{K,\Lambda}\bm w,
    \qquad
    g_t(\beta,\Lambda)=r_{\beta,t}a_t.
\]
Then
\[
    \Psi_T(\beta,\Lambda)=P_Tg(\beta,\Lambda)
    =T^{-1}\sum_{t=1}^Tg_t(\beta,\Lambda)
    =\innerT{\bm r_\beta}{\bm a}.
\]
Thus the sieve basis \(Q_{K,\Lambda}\) has already represented the nuisance
feasible directions: \(\operatorname{col}(Q_{K,\Lambda})\) is the tangent plane
for nuisance variation, and \(M_{K,\Lambda}\) keeps only the orthogonal score
component.

EL adds one more finite-dimensional convex problem. The primal feasible set is
the probability simplex cut by the affine moment hyperplane
\[
    \mathcal C_T(\beta,\Lambda)
    =\left\{\pi\in\R^T_+:
       \sum_{t=1}^T\pi_t=1,
       \quad \sum_{t=1}^T\pi_tg_t(\beta,\Lambda)=0
     \right\}.
\]
The empirical likelihood problem is equivalently
\begin{equation}
    \min_{\pi\in\mathcal C_T(\beta,\Lambda)}
    -\sum_{t=1}^T\log(T\pi_t).
    \label{eq:el-primal-convex}
\end{equation}
This is a convex barrier objective over a convex feasible set. Its Lagrangian is
\[
    \mathcal L(\pi,\eta,\lambda)
    =-\sum_{t=1}^T\log(T\pi_t)
      +\eta\left(\sum_{t=1}^T\pi_t-1\right)
      +T\lambda\sum_{t=1}^T\pi_tg_t.
\]
The KKT first-order condition gives
\[
    -\frac{1}{\pi_t}+\eta+T\lambda g_t=0.
\]
After normalizing by \(\sum_t\pi_t=1\), the EL weights are
\begin{equation}
    \widehat\pi_t(\lambda)
    =\frac{1}{T\{1+\lambda g_t(\beta,\Lambda)\}},
    \qquad
    1+\lambda g_t(\beta,\Lambda)>0.
    \label{eq:el-kkt-weights}
\end{equation}
The remaining KKT equation is the dual moment equation
\begin{equation}
    P_T\left[\frac{g_t(\beta,\Lambda)}
    {1+\lambda g_t(\beta,\Lambda)}\right]=0.
    \label{eq:el-dual-kkt}
\end{equation}
The log empirical likelihood ratio is therefore
\begin{equation}
    \ell_T(\beta,\Lambda)
    =2\sum_{t=1}^T\log\{1+\lambda_Tg_t(\beta,\Lambda)\}.
    \label{eq:el-ratio-kkt}
\end{equation}

The quadratic EL reduction is also mostly KKT algebra. If
\(\max_t|\lambda_Tg_t|=o_p(1)\), then expanding
\eqref{eq:el-dual-kkt} gives
\[
    0=P_T\left[\frac{g_t}{1+\lambda_Tg_t}\right]
      =P_Tg_t-\lambda_TP_Tg_t^2+R_{1T},
    \qquad
    R_{1T}=o_p(T^{-1/2}).
\]
Hence, with \(V_T=P_Tg_t^2\),
\[
    \lambda_T=\frac{P_Tg_t}{V_T}+o_p(T^{-1/2}).
\]
Expanding \eqref{eq:el-ratio-kkt} then yields
\begin{equation}
    \ell_T(\beta_0,\Lambda)
    =\frac{T\{P_Tg_t(\beta_0,\Lambda)\}^2}{V_T(\Lambda)}+o_p(1)
    =\frac{S_T(\beta_0,\Lambda)^2}{V_T(\Lambda)}+o_p(1).
    \label{eq:el-kkt-quadratic}
\end{equation}
The stochastic work is not to rediscover EL; it is only to verify the small size
conditions under which the KKT expansion is valid.

\section{Population Hilbert Space}

\subsection{Population object in \texorpdfstring{\(L^2(P)\)}{L2P}}

To express breaks in population geometry, include rescaled time \(S\in(0,1]\).
Let
\[
    O=(Y,X,W,S,U)
\]
be a population random element with law \(P\). The true regime map is
\[
    j_0(S)=j
    \quad\Longleftrightarrow\quad
    \tau_{j,0}<S\le\tau_{j+1,0}.
\]
The population model is the identity in \(\LtwoP\):
\begin{equation}
    Y=\beta_0X+\mu_0(S,W)+U,
    \qquad
    \mu_0(S,W)=\sum_{j=0}^{q_0}m_j(W)
    \ind\{\tau_{j,0}<S\le\tau_{j+1,0}\}.
    \label{eq:l2p-model}
\end{equation}

\subsection{Population nuisance space}

For a population partition \(\Lambda\), define
\[
    \N_\Lambda^P
    =\overline{\left\{
      \sum_{j=0}^{q} \ind\{S\in I_j(\Lambda)\}h_j(W):
      h_j\in L^2(P_W)
    \right\}}^{L^2(P)}.
\]
Let \(\Pi_\Lambda^P\) be the \(\LtwoP\)-orthogonal projection onto
\(\N_\Lambda^P\), and let
\[
    M_\Lambda^P=I-\Pi_\Lambda^P.
\]
For \(\beta\in\mathcal B\), define the population partial outcome
\[
    Y_\beta^\circ=Y-\beta X.
\]
The population profiled residual and score direction are
\[
    r_\beta^P(\Lambda)=M_\Lambda^PY_\beta^\circ
    =M_\Lambda^P(Y-\beta X),
    \qquad
    v^P(\Lambda)=M_\Lambda^Pw.
\]
The population moment is
\begin{equation}
    \Psi(\beta,\Lambda)
    =\innerP{r_\beta^P(\Lambda)}{v^P(\Lambda)}.
    \label{eq:population-moment}
\end{equation}

\section{Population Identification}

At the true population partition \(\Lambda_0\), abbreviate
\[
    \N_0^P=\N_{\Lambda_0}^P,
    \qquad
    \Pi_0^P=\Pi_{\Lambda_0}^P,
    \qquad
    M_0^P=M_{\Lambda_0}^P,
    \qquad
    v_0^P=v^P(\Lambda_0)=M_0^Pw.
\]
Because the true nuisance is regime-specific with the true population partition,
\[
    \mu_0(S,W)\in\N_0^P.
\]
Therefore the true-partition population residual-maker removes it:
\begin{equation}
    M_0^P\mu_0=0.
    \label{eq:population-nuisance-cancellation}
\end{equation}
This is the \(L^2(P)\) analog of the finite-sample statement that the correct
break-aware nuisance space absorbs the nuisance component.

Using the population model \eqref{eq:l2p-model}, the population partial outcome
at a generic \(\beta\) is
\[
\begin{aligned}
    Y_\beta^\circ
    &=Y-\beta X \\
    &=\mu_0(S,W)+U-(\beta-\beta_0)X.
\end{aligned}
\]
Projecting onto \((\N_0^P)^{\perp_P}\) gives
\begin{equation}
\begin{aligned}
    r_\beta^P(\Lambda_0)
    &=M_0^PY_\beta^\circ \\
    &=M_0^P\mu_0+M_0^PU-(\beta-\beta_0)M_0^PX \\
    &=M_0^PU-(\beta-\beta_0)M_0^PX,
\end{aligned}
\label{eq:population-profiled-residual-identity}
\end{equation}
because \(M_0^P\mu_0=0\). Taking the inner product with the true population
score direction \(v_0^P\) gives
\begin{equation}
    \Psi(\beta,\Lambda_0)
    =\innerP{M_0^PU}{v_0^P}
     -(\beta-\beta_0)\innerP{M_0^PX}{v_0^P}.
    \label{eq:population-identification}
\end{equation}
This derivation is algebraic. The two population assumptions enter only after
\eqref{eq:population-identification} is obtained.

The exogeneity condition is
\[
    \innerP{M_0^PU}{v_0^P}=0.
\]
The relevance condition is
\[
    \Gamma=\innerP{M_0^PX}{v_0^P}\ne0.
\]
Geometrically this is a transversality condition. After quotienting out the
nuisance tangent space, the parametric direction is not invisible to the chosen
normal score direction:
\[
    M_0^PX\not\perp_P M_0^Pw.
\]
Equivalently, the \(\beta\)-line crosses the population moment hyperplane
non-tangentially. If \(M_0^PX\) were orthogonal to every usable residualized
score direction, the parametric direction would be swallowed by the nuisance
space and \(\beta\) would not be identified by this moment.

Under these two conditions,
\[
    \Psi(\beta,\Lambda_0)=-(\beta-\beta_0)\Gamma.
\]
Thus
\begin{equation}
    \Psi(\beta,\Lambda_0)=0
    \quad\Longleftrightarrow\quad
    \beta=\beta_0.
    \label{eq:unique-identification}
\end{equation}
\part{Small Stochastic Bounds}

\section{Small Stochastic Bounds}

Part I has already fixed the objects. For every candidate partition \(\Lambda\),
write
\[
    \bm r_\beta(\Lambda)=\bm M_{K,\Lambda}(\bm y-\beta\bm x),
    \qquad
    \bm a(\Lambda)=\bm M_{K,\Lambda}\bm w,
\]
and
\[
    g_t(\beta,\Lambda)=r_{\beta,t}(\Lambda)a_t(\Lambda),
    \qquad
    \Psi_T(\beta,\Lambda)=P_Tg(\beta,\Lambda)
    =\innerT{\bm r_\beta(\Lambda)}{\bm a(\Lambda)}.
\]
At the true partition, abbreviate
\[
    \bm M_0=\bm M_{K,\Lambda_0},
    \qquad
    \bm v_0=\bm M_0\bm w,
    \qquad
    \bm y_\beta^\circ=\bm y-\beta\bm x.
\]
For the estimated partition, write
\[
    \widehat{\bm M}=\bm M_{K,\widehat\Lambda},
    \qquad
    \widehat{\bm v}=\widehat{\bm M}\bm w.
\]
The variance objects are
\[
    V_T(\Lambda)=P_Tg_t(\beta_0,\Lambda)^2,
    \qquad
    V_T^0=V_T(\Lambda_0),
    \qquad
    \widehat V_T=V_T(\widehat\Lambda).
\]

The stochastic part proves the following targets one by one:
\begin{equation}
    \sup_{\beta\in\mathcal B}
    |\Psi_T(\beta,\widehat\Lambda)-\Psi(\beta,\Lambda_0)|=\op(1),
    \label{eq:uniform-bridge}
\end{equation}
\[
    \sqrt T\Psi_T(\beta_0,\widehat\Lambda)\Rightarrow N(0,V),
    \qquad
    \widehat V_T\to_p V>0,
\]
and the local empirical-likelihood conditions that make the KKT expansion
\eqref{eq:el-kkt-quadratic} valid. None of these statements changes the
geometric objects. They only control their sampling error.

\subsection{Standing stochastic conditions}

Let \(\mathcal L_T\) denote the admissible set of candidate partitions searched
by the estimator. For local break-date arguments, let
\(\mathcal N_T(r_T)\subset\mathcal L_T\) denote the event/set of partitions whose
break dates are within the localization radius \(r_T\) of \(\Lambda_0\). The
following assumptions are the stochastic input for Part II. They are stated in
the workbook notation so each proof below can be read directly.

\begin{assumption}[Empirical-to-population inner-product bridge]
\label{ass:gram-stability}
Uniformly over \(\Lambda\in\mathcal L_T\), the empirical Gram matrices of the
block sieve design are stable:
\[
    \sup_{\Lambda\in\mathcal L_T}
    \left\|T^{-1}\bm Q_{K,\Lambda}'\bm Q_{K,\Lambda}-G_Q(\Lambda)\right\|_{op}
    =\op(1),
\]
where
\[
    0<c\le \lambda_{\min}\{G_Q(\Lambda)\}
    \le \lambda_{\max}\{G_Q(\Lambda)\}\le C<\infty
    \qquad \text{uniformly over }\Lambda\in\mathcal L_T.
\]
The same uniform empirical laws hold for the projected cross moments involving
\(\bm Q_{K,\Lambda}\), \(\bm x\), \(\bm w\), and \(\bm u\). Equivalently, for
the finite sieve and break classes used here,
\[
    \sup_{f,g\in\mathcal F_{K,T}}
    |\innerT{\bm f_T}{\bm g_T}-\innerP{f}{g}|=\op(1).
\]
This is the bridge between \(L^2(P_T)\) and \(L^2(P)\). The spaces are not being
identified as the same space; only the relevant inner products converge.
\end{assumption}

\begin{assumption}[Score-negligible sieve remainder]
\label{ass:score-negligible-sieve-remainder}
For the true nuisance vector \(\bm\mu_0\), the true-partition finite-sieve
remainder
\[
    \bm r_{\mu,T}=\bm M_0\bm\mu_0
\]
satisfies
\[
    \|\bm r_{\mu,T}\|_T=\op(1),
    \qquad
    \sqrt T\,|\innerT{\bm r_{\mu,T}}{\bm v_0}|=\op(1),
    \qquad
    P_T\{r_{\mu,T,t}^2v_{0,t}^2\}=\op(1).
\]
The first bound is the consistency-scale approximation bound. The second is the
score-scale bias bound. The third removes the sieve remainder from the variance.
\end{assumption}

\begin{assumption}[Estimated-partition primitive closure]
\label{ass:directional-break-replacement}
The RSS break estimator localizes at radius \(r_T\):
\[
    \Pr\{\widehat\Lambda\in\mathcal N_T(r_T)\}\to1.
\]
On \(\mathcal N_T(r_T)\), local perturbations of the projected score and second
moment are negligible:
\[
    \sup_{\Lambda\in\mathcal N_T(r_T)}
    \sqrt T\,|\Psi_T(\beta_0,\Lambda)-\Psi_T(\beta_0,\Lambda_0)|=\op(1),
\]
\[
    \sup_{\Lambda\in\mathcal N_T(r_T)}
    |V_T(\Lambda)-V_T^0|=\op(1),
\]
and, at the unscaled consistency level,
\[
    \sup_{\Lambda\in\mathcal N_T(r_T)}\sup_{\beta\in\mathcal B}
    |\Psi_T(\beta,\Lambda)-\Psi_T(\beta,\Lambda_0)|=\op(1).
\]
A sufficient local perturbation route is
\[
    \sup_{\Lambda\in\mathcal N_T(r_T)}
    T^{-1/2}|\bm r_{\mu,T}'(\bm M_{K,\Lambda}-\bm M_0)\bm w|=\op(1),
\]
\[
    \sup_{\Lambda\in\mathcal N_T(r_T)}
    T^{-1/2}|\bm u'(\bm M_{K,\Lambda}-\bm M_0)\bm w|=\op(1),
\]
with the variance replacement below. For a one-break local window, a standard
sufficient rate is
\[
    \zeta_K\sqrt{\frac{Kr_T}{T}}\to0,
\]
where \(\zeta_K\) controls the sieve-envelope size.
\end{assumption}

\begin{assumption}[RSS concentration and penalty rates]
\label{ass:rss-localization-penalty}
Let \(\Delta_\Lambda=\bm M_{K,\Lambda}-\bm M_0\). Outside the local set
\(\mathcal N_T(r_T)\), the deterministic RSS drift dominates stochastic RSS
fluctuation:
\[
    \inf_{\Lambda\in\mathcal L_T\setminus\mathcal N_T(r_T)}
    T^{-1}\bm\mu_0'\Delta_\Lambda\bm\mu_0\ge c_T,
\]
\[
    \sup_{\Lambda\in\mathcal L_T\setminus\mathcal N_T(r_T)}
    T^{-1}|2\bm\mu_0'\Delta_\Lambda\bm u+\bm u'\Delta_\Lambda\bm u|=o_p(c_T).
\]
One primitive route for the stochastic bound is conditionally sub-Gaussian
innovations plus sieve entropy/concentration over the searched RSS class.

For unknown order, the underfitted deterministic loss and overfitted stochastic
gain satisfy
\[
    \inf_{q<q_0}\{RSS_K(q)-RSS_K(q_0)\}\ge cT\Delta_T^2+o_p(T\Delta_T^2),
\]
\[
    \sup_{q>q_0}\{RSS_K(q_0)-RSS_K(q)\}=\Op(K\log T),
\]
and the penalty obeys
\[
    K\log T=o\{\operatorname{pen}_T(q,K)\},
    \qquad
    \operatorname{pen}_T(q,K)=o(T\Delta_T^2)
\]
for every fixed nonzero order difference.
\end{assumption}

\begin{assumption}[Projected LLN and CLT at the true partition]
\label{ass:projected-lln-clt}
The true-partition projected empirical moments obey
\[
    \innerT{\bm M_0\bm u}{\bm v_0}
    -\innerP{M_0^PU}{v_0^P}=\op(1),
\]
\[
    \innerT{\bm M_0\bm x}{\bm v_0}
    -\innerP{M_0^PX}{v_0^P}=\op(1),
\]
and
\[
    T^{-1/2}\bm u'\bm v_0\Rightarrow N(0,V),
    \qquad 0<V<\infty.
\]
\end{assumption}

\begin{assumption}[Second and higher moment control]
\label{ass:moment-control}
With \(\tilde g_t^0=(\bm M_0\bm u)_t v_{0,t}\),
\[
    P_T(\tilde g_t^0)^2\to_p V,
    \qquad
    P_T|g_t(\beta_0,\Lambda_0)|^3=\Op(1),
    \qquad
    P_T|g_t(\beta_0,\widehat\Lambda)|^3=\Op(1),
\]
and
\[
    \max_{t\le T}|g_t(\beta_0,\Lambda_0)|/\sqrt T=\op(1),
    \qquad
    \max_{t\le T}|g_t(\beta_0,\widehat\Lambda)|/\sqrt T=\op(1).
\]
\end{assumption}

\begin{assumption}[Scalar convex-hull condition]
\label{ass:convex-hull}
With probability tending to one, both scalar samples
\(\{g_t(\beta_0,\Lambda_0):1\le t\le T\}\) and
\(\{g_t(\beta_0,\widehat\Lambda):1\le t\le T\}\) contain at least one strictly
positive value and at least one strictly negative value.
\end{assumption}

\subsection{Primitive stochastic assumptions sufficient for Part II}

The high-level assumptions above are implied by familiar primitive bounds. Let
\[
    v^*_{t,0}=v_{0,t}=(\bm M_0\bm w)_t,
    \qquad
    \widehat v^*_{t,0}=(\widehat{\bm M}\bm w)_t.
\]
A useful sufficient primitive package is:
\begin{enumerate}[leftmargin=2em]
\item \textbf{Envelope control.}
      \[
          \max_{t\le T}|v^*_{t,0}|=\Op(1+\zeta_K),
          \qquad
          \max_{t\le T}|\widehat v^*_{t,0}|=\Op(1+\zeta_K).
      \]
\item \textbf{Variance replacement.} If
      \(\sigma_t^2=\E(u_t^2\mid\mathcal F_{t-1})\), then
      \[
          T^{-1}\sum_{t=1}^T
          \sigma_t^2\{(\widehat v^*_{t,0})^2-(v^*_{t,0})^2\}=\op(1).
      \]
\item \textbf{Local sieve-bias stability.}
      \[
          \sup_{\Lambda\in\mathcal N_T(r_T)}
          T^{-1/2}|\bm r_{\mu,T}'(\bm M_{K,\Lambda}-\bm M_0)\bm w|=\op(1).
      \]
\item \textbf{Local noise perturbation.}
      \[
          \sup_{\Lambda\in\mathcal N_T(r_T)}
          T^{-1/2}|\bm u'(\bm M_{K,\Lambda}-\bm M_0)\bm w|=\op(1),
      \]
      for example under the rate \(\zeta_K\sqrt{Kr_T/T}\to0\) in the one-break
      local perturbation calculation.
\item \textbf{RSS concentration.} The errors have enough tail control, for
      example conditional sub-Gaussian tails, so the searched RSS class satisfies
      the concentration bounds in Assumption~\ref{ass:rss-localization-penalty}.
\item \textbf{Penalty separation.} The unknown-order penalty dominates searched
      overfit gains but is smaller than the underfit deterministic loss:
      \[
          K\log T=o\{\operatorname{pen}_T(q,K)\},
          \qquad
          \operatorname{pen}_T(q,K)=o(T\Delta_T^2).
      \]
\end{enumerate}
These primitive conditions are not new geometry. They are the stochastic closure
that makes the geometric KKT construction usable with estimated breaks.
\subsection{Part II primitive proof-claim closure}

This subsection records the detailed primitive route behind the high-level
Part II assumptions. It is not a second model. It is the proof claim that the
RSS localization, projected score CLT, variance replacement, max-score bound,
and EL KKT expansion can all be obtained from standard local projection and
concentration bounds.

Write
\[
    d_K=(q_0+1)K,\qquad b_K=1+\zeta_K,
    \qquad \bm\Pi_0=I-\bm M_0.
\]
At the true partition let \(\bm Q_0=\bm Q_{K,\Lambda_0}\). Assume the true
nuisance has the finite-sieve decomposition
\[
    \bm\mu_0=\bm Q_0c_0+\bm r_{\mu,T}.
\]

\begin{lemma}[Scalar EL KKT proof claim]
\label{lem:part2-scalar-el-proof-claim}
Let \(Z_{t,T}\) be any scalar triangular moment array, and define
\[
    A_T=T^{-1/2}\sum_{t=1}^T Z_{t,T},\qquad
    B_T=T^{-1}\sum_{t=1}^T Z_{t,T}^2,\qquad
    m_T=\max_{t\le T}|Z_{t,T}|.
\]
If \(A_T=\Op(1)\), \(B_T\ge c+\op(1)\), \(m_T=\op(\sqrt T)\), and the scalar
convex hull contains zero in its interior with probability tending to one, then
the EL dual root \(\lambda_T\) satisfies
\[
    \lambda_T=\frac{T^{-1}\sum_t Z_{t,T}}{T^{-1}\sum_t Z_{t,T}^2}
    +\op(T^{-1/2}),
\]
and
\[
    2\sum_{t=1}^T\log(1+\lambda_TZ_{t,T})
    =\frac{A_T^2}{B_T}+\op(1).
\]
\end{lemma}

\begin{proof}
The KKT equation is
\[
    0=\sum_t\frac{Z_{t,T}}{1+\lambda_TZ_{t,T}},
    \qquad 1+\lambda_TZ_{t,T}>0.
\]
With \(S_1=\sum_tZ_{t,T}\) and \(S_2=\sum_tZ_{t,T}^2\), the exact identity
\[
    S_1=\lambda_T\sum_t\frac{Z_{t,T}^2}{1+\lambda_TZ_{t,T}}
\]
implies \(|\lambda_T|=\Op(T^{-1/2})\) because
\(S_1=\Op(\sqrt T)\), \(S_2\ge cT+\op(T)\), and
\(|S_1|m_T=\op(T)\). Hence \(|\lambda_T|m_T=\op(1)\). Expanding the KKT
equation around zero gives
\[
    0=S_1-\lambda_TS_2+\op(\sqrt T),
    \qquad
    \lambda_T=S_1/S_2+\op(T^{-1/2}).
\]
Taylor expansion of the log ratio gives
\[
    2\sum_t\log(1+\lambda_TZ_{t,T})
    =2\lambda_TS_1-\lambda_T^2S_2+\op(1)
    =S_1^2/S_2+\op(1).
\]
Since \(S_1^2/S_2=A_T^2/B_T\), the claim follows.
\end{proof}

\begin{lemma}[True-partition feasible-to-oracle proof claim]
\label{lem:part2-feasible-oracle-proof-claim}
At \((\beta_0,\Lambda_0)\),
\[
    \sqrt T\Psi_T(\beta_0,\Lambda_0)
    =T^{-1/2}\bm u'\bm v_0
     +T^{-1/2}\bm r_{\mu,T}'\bm v_0.
\]
If \(T^{-1/2}\bm r_{\mu,T}'\bm v_0=\op(1)\), then the score numerator is
asymptotically the oracle numerator \(T^{-1/2}\bm u'\bm v_0\). Moreover, if
\[
    \|(\bm M_0\bm r_{\mu,T}-\bm\Pi_0\bm u)\odot\bm v_0\|_T=\op(1),
    \qquad
    \|(\bm M_0\bm r_{\mu,T}-\bm\Pi_0\bm u)\odot\bm v_0\|_\infty=\op(\sqrt T),
\]
then
\[
    V_T^0=T^{-1}\sum_{t=1}^T u_t^2v_{0,t}^2+\op(1),
    \qquad
    \max_t|g_t(\beta_0,\Lambda_0)|
    =\max_t|u_tv_{0,t}|+\op(\sqrt T).
\]
\end{lemma}

\begin{proof}
Since \(\bm M_0\bm Q_0=0\),
\[
    \bm M_0(\bm y-\beta_0\bm x)
    =\bm M_0(\bm\mu_0+\bm u)
    =\bm r_{\mu,T}+\bm M_0\bm u.
\]
Because \(\bm v_0=\bm M_0\bm w\) and \(\bm M_0\) is symmetric and idempotent,
\[
    (\bm M_0\bm u)'\bm v_0=\bm u'\bm v_0.
\]
This proves the numerator identity.

For the variance and max bounds, write
\[
    \bm M_0\bm u=\bm u-\bm\Pi_0\bm u
\]
and therefore
\[
    g_t(\beta_0,\Lambda_0)
    =u_tv_{0,t}+d_{t,0},
    \qquad
    \bm d_0=(\bm M_0\bm r_{\mu,T}-\bm\Pi_0\bm u)\odot\bm v_0.
\]
The \(L^2(P_T)\) bound on \(\bm d_0\) gives the second-moment replacement by
Cauchy-Schwarz, using \(T^{-1}\sum_tu_t^2v_{0,t}^2=\Op(1)\). The sup-norm
bound gives the max-score replacement.
\end{proof}

The preceding display is the place where primitive projection rates enter. A
standard sufficient package is
\[
    \|\bm M_0\bm r_{\mu,T}\|_T=\Op(a_K),\qquad
    \|\bm\Pi_0\bm u\|_T=\Op\!\left(\sqrt{d_K/T}\right),
\]
\[
    \|\bm M_0\bm r_{\mu,T}\|_\infty=\Op(a_Kb_K),\qquad
    \|\bm\Pi_0\bm u\|_\infty
    =\Op\!\left(\zeta_K\sqrt{d_K/T}\right),
\]
together with \(\|\bm v_0\|_\infty=\Op(b_K)\). These imply
\[
    \|\bm d_0\|_T
    =\Op\!\left(a_Kb_K+b_K\sqrt{d_K/T}\right)=\op(1),
\]
and
\[
    \|\bm d_0\|_\infty
    =\Op\!\left(a_Kb_K^2+\zeta_Kb_K\sqrt{d_K/T}\right)=\op(\sqrt T).
\]

\begin{lemma}[Oracle score proof claim]
\label{lem:part2-oracle-score-proof-claim}
Suppose either \(\bm v_0\) is admissible for the martingale array, or there is a
predictable approximation \(v_{t,0}^{\mathrm{pr}}\) satisfying
\[
    T^{-1/2}\sum_t u_t(v_{0,t}-v_{t,0}^{\mathrm{pr}})=\op(1),
\]
\[
    T^{-1}\sum_t u_t^2\{v_{0,t}^2-(v_{t,0}^{\mathrm{pr}})^2\}=\op(1),
    \qquad
    \max_t|v_{t,0}^{\mathrm{pr}}|=\Op(b_K).
\]
If \(b_K^{2+\delta_0}T^{-\delta_0/2}\to0\) for some
\(\delta_0>0\), then
\[
    T^{-1/2}\sum_tu_tv_{0,t}\Rightarrow N(0,V),
\]
\[
    T^{-1}\sum_tu_t^2v_{0,t}^2\to_p V,
    \qquad
    \max_t|u_tv_{0,t}|=\op(\sqrt T).
\]
\end{lemma}

\begin{proof}
Apply the martingale triangular-array CLT to
\(\xi_{t,T}=T^{-1/2}u_tv_{t,0}^{\mathrm{pr}}\). The conditional variance is
\[
    T^{-1}\sum_t\sigma_t^2(v_{t,0}^{\mathrm{pr}})^2\to_p V,
\]
and the conditional Lyapunov term is bounded by
\[
    CT^{-(1+\delta_0/2)}
    \sum_t|v_{t,0}^{\mathrm{pr}}|^{2+\delta_0}
    \le Cb_K^{2+\delta_0}T^{-\delta_0/2}=\op(1).
\]
The predictable-approximation assumption transfers the CLT back to \(\bm v_0\).
The variance LLN follows by applying the same martingale moment argument to
\((u_t^2-\sigma_t^2)(v_{t,0}^{\mathrm{pr}})^2\), and the max bound follows from
\[
    \Pp\{\max_t|u_tv_{t,0}^{\mathrm{pr}}|>\varepsilon\sqrt T\}
    \le C_\varepsilon b_K^{2+\delta_0}T^{-\delta_0/2}+o(1).
\]
\end{proof}

\begin{lemma}[Local estimated-partition bridge proof claim]
\label{lem:part2-local-bridge-proof-claim}
For one break, let \(\Lambda_0=(k_0)\), \(\widetilde\Lambda=(\widetilde k)\),
and \(h=|\widetilde k-k_0|=\Op(r_T)\). Under local Gram stability,
\[
    T^{-1/2}\bm\mu_0'(\bm M_{K,\widetilde k}-\bm M_0)\bm w
    =\Op\!\left(\frac{b_Kr_T}{\sqrt T}\right)+\op(1),
\]
and
\[
    T^{-1/2}\bm u'(\bm M_{K,\widetilde k}-\bm M_0)\bm w
    =\Op\!\left(\zeta_K\sqrt{\frac{Kr_T}{T}}\right)+\op(1).
\]
If
\[
    \frac{b_Kr_T}{\sqrt T}\to0,
    \qquad
    \zeta_K\sqrt{\frac{Kr_T}{T}}\to0,
\]
then
\[
    \sqrt T\{\Psi_T(\beta_0,\widetilde\Lambda)
    -\Psi_T(\beta_0,\Lambda_0)\}=\op(1).
\]
If also \(\zeta_K^2Kr_T/T\to0\), then
\[
    V_T(\widetilde\Lambda)-V_T(\Lambda_0)=\op(1).
\]
\end{lemma}

\begin{proof}
Take \(\widetilde k=k_0+h\); the other side is symmetric. The only observations
whose regime label changes are in \(H=\{k_0+1,\ldots,k_0+h\}\). Relative to the
candidate partition, the true nuisance can be written as
\[
    \bm\mu_0=\bm Q_{K,\widetilde k}c^{(\widetilde k)}
      +\bm d_H+\bm r_{\mu,T},
\]
where \(\bm d_H\) is supported on \(H\) and \(\|\bm d_H\|_\infty=\Op(1)\). Since
\(\bm M_{K,\widetilde k}\bm Q_{K,\widetilde k}=0\) and
\(\bm M_0\bm Q_0=0\),
\[
    \bm\mu_0'(\bm M_{K,\widetilde k}-\bm M_0)\bm w
    =\bm d_H'\bm M_{K,\widetilde k}\bm w
      +\bm r_{\mu,T}'(\bm M_{K,\widetilde k}-\bm M_0)\bm w.
\]
The first term is bounded by \(\|\bm d_H\|_1\|\bm M_{K,\widetilde k}\bm w\|_\infty
=\Op(hb_K)\), and the second term is the local sieve-bias stability term.
This proves the nuisance perturbation bound.

For the stochastic perturbation, write the residualized weights on the three
blocks \(A=\{1,\ldots,k_0\}\), \(H\), and \(B=\{k_0+h+1,\ldots,T\}\). The
least-squares coefficient update on a long segment has the form
\[
    \widehat\gamma_{A\cup H}-\widehat\gamma_A
    =(P_{A\cup H}'P_{A\cup H})^{-1}
      P_H'(w_H-P_H\widehat\gamma_A),
\]
so local Gram stability gives
\[
    \|\widehat\gamma_{A\cup H}-\widehat\gamma_A\|
    =\Op(h\zeta_Kb_K/T),
\]
and similarly on the right side. After cancellation, the remaining stochastic
terms are controlled by
\[
    \|P_A'u_A\|+\|P_B'u_B\|=\Op(\sqrt{TK}),
    \qquad
    \|P_H'u_H\|=\Op(\zeta_K\sqrt h).
\]
Multiplying by \(T^{-1/2}\) gives the displayed
\(\Op\{\zeta_K\sqrt{Kr_T/T}\}\) bound. The denominator replacement follows by
the same local boundary argument: \(v_{\widetilde k}-v_{k_0}\) is order
\(\Op(b_K)\) only on the \(h\)-point boundary and order
\(\Op(h\zeta_K^2b_K/T)\) away from it, so
\[
    |V_T(\widetilde\Lambda)-V_T(\Lambda_0)|
    =\Op\!\left(\frac{b_K^2r_T}{T}
      +\frac{\zeta_K^2Kr_T}{T}\right)+\op(1).
\]
The stated rates make this \(\op(1)\). A fixed finite number of breaks is
handled by summing the same boundary calculation over the finitely many local
neighborhoods.
\end{proof}

\begin{proposition}[Localization and order-selection proof claim]
\label{prop:part2-localization-order-proof-claim}
In the one-break calculation, suppose
\[
    \Delta RSS_K(k)
    =RSS_K(\beta_0,k)-RSS_K(\beta_0,k_0)
\]
satisfies, uniformly for \(h=|k-k_0|\),
\[
    E\{\Delta RSS_K(k)\mid D_T\}
    \ge c h\Delta_T^2-o_p(h\Delta_T^2),
\]
and
\[
    \left|\Delta RSS_K(k)-E\{\Delta RSS_K(k)\mid D_T\}\right|
    =\Op(\sqrt{hK\log T})+\Op(K\log T).
\]
Then with \(R_T=K\log T\) and
\[
    r_T=R_T\max\{\Delta_T^{-2},\Delta_T^{-4}\},
\]
the known-order break estimator satisfies
\[
    |\widehat k-k_0|=\Op(r_T).
\]
If \(\Delta_T\) is bounded away from zero, this reduces to
\(|\widehat k-k_0|=\Op(K\log T)\). With fixed \(\bar q\), if the overfit
penalty increment dominates \(K\log T\) and the underfit penalty saving is
\(o(T\Delta_T^2)\), then \(\Pp(\widehat q=q_0)\to1\).
\end{proposition}

\begin{proof}
The RSS identity from the geometric part gives
\[
    \Delta RSS_K(k)
    =\bm\mu_0'A_k\bm\mu_0
      +2\bm\mu_0'A_k\bm u+\bm u'A_k\bm u,
    \qquad A_k=\bm M_{K,k}-\bm M_{K,k_0}.
\]
Conditional centering and equal-rank residual-makers give
\(E(\bm u'A_k\bm u\mid D_T)=\sigma^2\operatorname{tr}(A_k)=0\), so the
conditional mean is the deterministic drift. The displayed stochastic
fluctuation bound comes from a sub-Gaussian bound for the linear term and a
Hanson-Wright bound for the rank-\(O(K)\) quadratic term, followed by a union
bound over searched break dates.

For \(h\ge Mr_T\),
\[
    \frac{R_T}{h\Delta_T^2}\le M^{-1},
    \qquad
    \frac{\sqrt{hR_T}}{h\Delta_T^2}\le M^{-1/2}.
\]
Choosing \(M\) large makes the stochastic part smaller than the positive drift,
so \(\Delta RSS_K(k)>0\) uniformly outside the \(Mr_T\) window with probability
tending to one. Since \(\widehat k\) minimizes RSS and \(k_0\) is admissible,
\(|\widehat k-k_0|=\Op(r_T)\).

For unknown order, underfitting loses at least \(cT\Delta_T^2+o_p(T\Delta_T^2)\)
in deterministic RSS, while the penalty saving is \(o(T\Delta_T^2)\). Overfitting
can improve RSS only by \(\Op(K\log T)\), while the penalty increase is larger
than \(K\log T\). A finite union over \(q\le\bar q\) gives
\(\Pp(\widehat q=q_0)\to1\).
\end{proof}
\subsection{RSS localization and order selection}

\begin{proposition}[RSS localization]
\label{prop:part2-rss-localization}
Under Assumption~\ref{ass:rss-localization-penalty}, the known-order RSS
estimator satisfies
\[
    \Pr\{\widehat\Lambda\in\mathcal N_T(r_T)\}\to1.
\]
\end{proposition}

\begin{proof}
For any candidate \(\Lambda\), Section~7 gives the exact decomposition
\[
\begin{aligned}
    RSS_K(\beta_0,\Lambda)-RSS_K(\beta_0,\Lambda_0)
    &=\bm\mu_0'\Delta_\Lambda\bm\mu_0
      +2\bm\mu_0'\Delta_\Lambda\bm u
      +\bm u'\Delta_\Lambda\bm u,
\end{aligned}
\]
where \(\Delta_\Lambda=\bm M_{K,\Lambda}-\bm M_0\). On
\(\mathcal L_T\setminus\mathcal N_T(r_T)\), Assumption~\ref{ass:rss-localization-penalty}
gives
\[
    T^{-1}\bm\mu_0'\Delta_\Lambda\bm\mu_0\ge c_T
\]
uniformly, while
\[
    T^{-1}|2\bm\mu_0'\Delta_\Lambda\bm u+\bm u'\Delta_\Lambda\bm u|=o_p(c_T)
\]
uniformly. Therefore
\[
    \inf_{\Lambda\in\mathcal L_T\setminus\mathcal N_T(r_T)}
    \{RSS_K(\beta_0,\Lambda)-RSS_K(\beta_0,\Lambda_0)\}>0
\]
with probability tending to one. A minimizer over the known-order search set
cannot lie outside \(\mathcal N_T(r_T)\) on this event.
\end{proof}

\begin{proposition}[Unknown-order penalty selection]
\label{prop:part2-order-selection}
Under Assumption~\ref{ass:rss-localization-penalty}, the penalized RSS selector
satisfies
\[
    \Pr(\widehat q=q_0)\to1.
\]
\end{proposition}

\begin{proof}
For underfitting, Assumption~\ref{ass:rss-localization-penalty} gives
\[
    \inf_{q<q_0}\{RSS_K(q)-RSS_K(q_0)\}
    \ge cT\Delta_T^2+o_p(T\Delta_T^2).
\]
The penalty difference for a fixed lower order is
\(o(T\Delta_T^2)\), so the deterministic loss from omitting a true break
dominates any penalty saving. Thus underfitted orders are rejected with
probability tending to one.

For overfitting,
\[
    \sup_{q>q_0}\{RSS_K(q_0)-RSS_K(q)\}=\Op(K\log T).
\]
The penalty increment satisfies \(K\log T=o\{\operatorname{pen}_T(q,K)\}\), so
any spurious RSS gain is dominated by the penalty increase with probability
tending to one. Thus overfitted orders are rejected with probability tending to
one. Combining the underfit and overfit cases gives \(\Pr(\widehat q=q_0)\to1\).
\end{proof}
\subsection{Gram stability}

\begin{lemma}[Gram stability gives well-behaved projections]
\label{lem:part2-gram-stability}
Under Assumption~\ref{ass:gram-stability}, uniformly over
\(\Lambda\in\mathcal L_T\),
\[
    \lambda_{\min}\{T^{-1}\bm Q_{K,\Lambda}'\bm Q_{K,\Lambda}\}\ge c/2
\]
with probability tending to one. Hence the coordinate formula for
\(\Pi_{K,\Lambda,T}\) is stable on \(\mathcal L_T\), and the residual-maker
\(\bm M_{K,\Lambda}=I-\Pi_{K,\Lambda,T}\) remains an orthogonal contraction in
\(\LtwoT\):
\[
    \|\bm M_{K,\Lambda}\bm z\|_T\le \|\bm z\|_T
    \qquad \text{for every }\bm z\in\LtwoT.
\]
\end{lemma}

\begin{proof}
By Weyl's eigenvalue inequality,
\[
\begin{aligned}
\lambda_{\min}\{T^{-1}\bm Q_{K,\Lambda}'\bm Q_{K,\Lambda}\}
&\ge
\lambda_{\min}\{G_Q(\Lambda)\}
-\left\|T^{-1}\bm Q_{K,\Lambda}'\bm Q_{K,\Lambda}-G_Q(\Lambda)\right\|_{op}.
\end{aligned}
\]
Taking the infimum over \(\Lambda\in\mathcal L_T\), Assumption~\ref{ass:gram-stability}
puts the first term above \(c\) and the second term below \(c/2\) with
probability tending to one. This proves the eigenvalue bound.

On that event, the least-squares projection has the usual stable coordinate
form
\[
    \Pi_{K,\Lambda,T}
    =\bm Q_{K,\Lambda}
      (\bm Q_{K,\Lambda}'\bm Q_{K,\Lambda})^{-1}
      \bm Q_{K,\Lambda}'.
\]
Independently of coordinates, \(\Pi_{K,\Lambda,T}\) is the orthogonal projection
onto \(\N_{K,\Lambda,T}\). Therefore
\(\bm M_{K,\Lambda}=I-\Pi_{K,\Lambda,T}\) is also an orthogonal projection.
Orthogonal projections are contractions in Hilbert norm.
\end{proof}

\subsection{Empirical-to-population projection bridge}

\begin{lemma}[Inner-product bridge from \(L^2(P_T)\) to \(L^2(P)\)]
\label{lem:part2-inner-product-bridge}
Under Assumption~\ref{ass:gram-stability}, uniformly over the sieve and break
classes used in the projected score,
\[
    \innerT{\bm f_T}{\bm g_T}=\innerP{f}{g}+\op(1).
\]
In particular,
\[
    \sup_{\Lambda\in\mathcal L_T}
    \left\|T^{-1}\bm Q_{K,\Lambda}'\bm Q_{K,\Lambda}-G_Q(\Lambda)\right\|_{op}
    =\op(1),
\]
and the corresponding projected \((x,w,u)\) cross moments converge to their
population limits.
\end{lemma}

\begin{proof}
The first display is the uniform empirical law in
Assumption~\ref{ass:gram-stability}. Taking \(f\) and \(g\) to be two
coordinates of the block sieve row gives the displayed Gram convergence.
Taking one or both functions to be the projected \(x\), \(w\), or \(u\)
directions gives the corresponding cross-moment convergence. This is exactly
the bridge needed by the workbook: finite-sample projection geometry is written
in \(L^2(P_T)\), while identification is written in \(L^2(P)\), and the common
object is convergence of the relevant inner products.
\end{proof}
\subsection{Sieve remainder is small}

\begin{lemma}[True-partition sieve remainder is negligible]
\label{lem:part2-sieve-remainder}
Under Assumption~\ref{ass:score-negligible-sieve-remainder},
\[
    \innerT{\bm M_0\bm\mu_0}{\bm v_0}=\op(1),
    \qquad
    \sqrt T\innerT{\bm M_0\bm\mu_0}{\bm v_0}=\op(1),
\]
and
\[
    V_T^0
    =P_T\{(\bm M_0\bm u)_tv_{0,t}\}^2+\op(1).
\]
\end{lemma}

\begin{proof}
By definition \(\bm r_{\mu,T}=\bm M_0\bm\mu_0\). The first two displays are
exactly the consistency-scale and score-scale components of
Assumption~\ref{ass:score-negligible-sieve-remainder}.

At \(\beta_0\),
\[
    \bm M_0\bm y_{\beta_0}^\circ
    =\bm M_0(\bm\mu_0+\bm u)
    =\bm r_{\mu,T}+\bm M_0\bm u.
\]
Thus
\[
    g_t(\beta_0,\Lambda_0)
    =\{(\bm M_0\bm u)_t+r_{\mu,T,t}\}v_{0,t}.
\]
Let \(\tilde g_t^0=(\bm M_0\bm u)_tv_{0,t}\) and
\(d_t=r_{\mu,T,t}v_{0,t}\). Then
\[
    V_T^0-P_T(\tilde g_t^0)^2
    =2P_T\tilde g_t^0d_t+P_Td_t^2.
\]
The final term is \(\op(1)\) by Assumption~\ref{ass:score-negligible-sieve-remainder}.
The cross term is bounded by Cauchy-Schwarz:
\[
    |P_T\tilde g_t^0d_t|
    \le \{P_T(\tilde g_t^0)^2\}^{1/2}\{P_Td_t^2\}^{1/2}=\op(1),
\]
using the bounded second moment in Assumption~\ref{ass:moment-control}. Hence
\(V_T^0=P_T(\tilde g_t^0)^2+\op(1)\).
\end{proof}

\subsection{Break replacement}

\begin{lemma}[Estimated partition may replace the true partition]
\label{lem:part2-break-replacement}
Under Assumption~\ref{ass:directional-break-replacement},
\[
    \sup_{\beta\in\mathcal B}
    |\Psi_T(\beta,\widehat\Lambda)-\Psi_T(\beta,\Lambda_0)|=\op(1),
\]
\[
    \sqrt T\,
    |\Psi_T(\beta_0,\widehat\Lambda)-\Psi_T(\beta_0,\Lambda_0)|=\op(1),
    \qquad
    \widehat V_T-V_T^0=\op(1).
\]
\end{lemma}

\begin{proof}
By Assumption~\ref{ass:directional-break-replacement},
\(\Pr\{\widehat\Lambda\in\mathcal N_T(r_T)\}\to1\). On this event,
\(\widehat\Lambda\) is one of the partitions over which the three local
suprema in the same assumption are taken. Therefore
\[
    \sup_{\beta\in\mathcal B}
    |\Psi_T(\beta,\widehat\Lambda)-\Psi_T(\beta,\Lambda_0)|
    \le
    \sup_{\Lambda\in\mathcal N_T(r_T)}\sup_{\beta\in\mathcal B}
    |\Psi_T(\beta,\Lambda)-\Psi_T(\beta,\Lambda_0)|=\op(1),
\]
\[
    \sqrt T|\Psi_T(\beta_0,\widehat\Lambda)-\Psi_T(\beta_0,\Lambda_0)|
    \le
    \sup_{\Lambda\in\mathcal N_T(r_T)}
    \sqrt T|\Psi_T(\beta_0,\Lambda)-\Psi_T(\beta_0,\Lambda_0)|=\op(1),
\]
and
\[
    |\widehat V_T-V_T^0|
    \le
    \sup_{\Lambda\in\mathcal N_T(r_T)}|V_T(\Lambda)-V_T^0|=\op(1).
\]
The conclusion follows after adding the negligible probability of the complement
of the localization event.
\end{proof}

\subsection{Uniform moment bridge}

\begin{lemma}[Uniform projected moment bridge]
\label{lem:part2-uniform-bridge}
Under Assumptions~\ref{ass:score-negligible-sieve-remainder},
\ref{ass:directional-break-replacement}, and \ref{ass:projected-lln-clt},
\begin{equation}
    \sup_{\beta\in\mathcal B}
    |\Psi_T(\beta,\widehat\Lambda)-\Psi(\beta,\Lambda_0)|=\op(1).
    \tag{\ref{eq:uniform-bridge}}
\end{equation}
\end{lemma}

\begin{proof}
First prove the bridge at the true partition. Since
\(\bm y_\beta^\circ=\bm\mu_0+\bm u-(\beta-\beta_0)\bm x\),
\[
\begin{aligned}
    \Psi_T(\beta,\Lambda_0)
    &=\innerT{\bm M_0\bm y_\beta^\circ}{\bm v_0} \\
    &=\innerT{\bm M_0\bm\mu_0}{\bm v_0}
      +\innerT{\bm M_0\bm u}{\bm v_0}
      -(\beta-\beta_0)\innerT{\bm M_0\bm x}{\bm v_0}.
\end{aligned}
\]
The population moment from Part I is
\[
    \Psi(\beta,\Lambda_0)
    =\innerP{M_0^PU}{v_0^P}
     -(\beta-\beta_0)\innerP{M_0^PX}{v_0^P}.
\]
Subtracting gives
\[
\begin{aligned}
    \Psi_T(\beta,\Lambda_0)-\Psi(\beta,\Lambda_0)
    &=\innerT{\bm M_0\bm\mu_0}{\bm v_0} \\
    &\quad+
      \left\{\innerT{\bm M_0\bm u}{\bm v_0}
      -\innerP{M_0^PU}{v_0^P}\right\} \\
    &\quad-(\beta-\beta_0)
      \left\{\innerT{\bm M_0\bm x}{\bm v_0}
      -\innerP{M_0^PX}{v_0^P}\right\}.
\end{aligned}
\]
The first term is \(\op(1)\) by Lemma~\ref{lem:part2-sieve-remainder}. The next
two bracketed terms are \(\op(1)\) by Assumption~\ref{ass:projected-lln-clt}.
Because \(\mathcal B\) is fixed and bounded, the convergence is uniform in
\(\beta\in\mathcal B\):
\[
    \sup_{\beta\in\mathcal B}
    |\Psi_T(\beta,\Lambda_0)-\Psi(\beta,\Lambda_0)|=\op(1).
\]
Finally, Lemma~\ref{lem:part2-break-replacement} gives
\[
    \sup_{\beta\in\mathcal B}
    |\Psi_T(\beta,\widehat\Lambda)-\Psi_T(\beta,\Lambda_0)|=\op(1).
\]
The triangle inequality proves the stated bridge.
\end{proof}

\subsection{Score CLT}

\begin{lemma}[Projected score CLT]
\label{lem:part2-score-clt}
Under Assumptions~\ref{ass:score-negligible-sieve-remainder},
\ref{ass:directional-break-replacement}, and \ref{ass:projected-lln-clt},
\[
    \sqrt T\Psi_T(\beta_0,\widehat\Lambda)
    \Rightarrow N(0,V).
\]
\end{lemma}

\begin{proof}
At the true partition,
\[
\begin{aligned}
    \sqrt T\Psi_T(\beta_0,\Lambda_0)
    &=\sqrt T\innerT{\bm M_0\bm y_{\beta_0}^\circ}{\bm v_0} \\
    &=T^{-1/2}\bm u'\bm v_0
      +\sqrt T\innerT{\bm M_0\bm\mu_0}{\bm v_0}.
\end{aligned}
\]
The first equality is the definition of \(\Psi_T\). The second uses
\(\bm y_{\beta_0}^\circ=\bm\mu_0+\bm u\) and the symmetry/idempotence of
\(\bm M_0\), so \(\innerT{\bm M_0\bm u}{\bm v_0}=T^{-1}\bm u'\bm v_0\).
The approximation-bias term is \(\op(1)\) by
Assumption~\ref{ass:score-negligible-sieve-remainder}. The leading term
\(T^{-1/2}\bm u'\bm v_0\) converges to \(N(0,V)\) by
Assumption~\ref{ass:projected-lln-clt}. Therefore
\[
    \sqrt T\Psi_T(\beta_0,\Lambda_0)\Rightarrow N(0,V).
\]
Lemma~\ref{lem:part2-break-replacement} gives
\[
    \sqrt T\{
    \Psi_T(\beta_0,\widehat\Lambda)-\Psi_T(\beta_0,\Lambda_0)
    \}=\op(1).
\]
Slutsky's theorem gives the result for \(\widehat\Lambda\).
\end{proof}

\subsection{Variance consistency}

\begin{lemma}[Projected variance consistency]
\label{lem:part2-variance-consistency}
Under Assumptions~\ref{ass:score-negligible-sieve-remainder},
\ref{ass:directional-break-replacement}, and \ref{ass:moment-control},
\[
    \widehat V_T\to_p V>0.
\]
\end{lemma}

\begin{proof}
Lemma~\ref{lem:part2-sieve-remainder} gives
\[
    V_T^0=P_T(\tilde g_t^0)^2+\op(1),
    \qquad
    \tilde g_t^0=(\bm M_0\bm u)_tv_{0,t}.
\]
Assumption~\ref{ass:moment-control} gives \(P_T(\tilde g_t^0)^2\to_p V\).
Therefore \(V_T^0\to_p V\). Lemma~\ref{lem:part2-break-replacement} gives
\(\widehat V_T-V_T^0=\op(1)\), so \(\widehat V_T\to_p V\). The lower bound
\(V>0\) is part of Assumption~\ref{ass:projected-lln-clt}; it rules out a
degenerate residual score direction.
\end{proof}

\subsection{Max-score and convex hull}

\begin{lemma}[EL local feasibility and max-score]
\label{lem:part2-el-feasibility}
Under Assumptions~\ref{ass:moment-control} and \ref{ass:convex-hull},
\[
    \max_{t\le T}|g_t(\beta_0,\Lambda_0)|/\sqrt T=\op(1),
    \qquad
    \max_{t\le T}|g_t(\beta_0,\widehat\Lambda)|/\sqrt T=\op(1),
\]
and the origin is in the interior of the scalar convex hull for both
\(\{g_t(\beta_0,\Lambda_0):1\le t\le T\}\) and
\(\{g_t(\beta_0,\widehat\Lambda):1\le t\le T\}\) with probability tending to
one.
\end{lemma}

\begin{proof}
The two max-score bounds are displayed in Assumption~\ref{ass:moment-control}.
For scalar moments, the convex hull of the observed values is the interval from
the sample minimum to the sample maximum. The origin is in its interior exactly
when the minimum is strictly negative and the maximum is strictly positive.
Assumption~\ref{ass:convex-hull} states that this sign condition holds with
probability tending to one for both the known-partition and estimated-partition
moment arrays.
\end{proof}

\subsection{EL KKT local expansion}

\begin{lemma}[Local EL KKT expansion]
\label{lem:part2-el-local-expansion}
Under Assumptions~\ref{ass:score-negligible-sieve-remainder}, \ref{ass:directional-break-replacement}, \ref{ass:projected-lln-clt}, \ref{ass:moment-control}, and \ref{ass:convex-hull},
with \(g_t=g_t(\beta_0,\widehat\Lambda)\),
\[
    \lambda_T=\frac{P_Tg_t}{P_Tg_t^2}+\op(T^{-1/2}),
    \qquad
    \max_{t\le T}|\lambda_Tg_t|=\op(1),
\]
and
\[
    \ell_T(\beta_0,\widehat\Lambda)
    =\frac{T\{\Psi_T(\beta_0,\widehat\Lambda)\}^2}{\widehat V_T}
     +\op(1).
\]
\end{lemma}

\begin{proof}
By Lemma~\ref{lem:part2-el-feasibility}, the scalar convex-hull condition holds
with probability tending to one, so the EL dual root \(\lambda_T\) exists in the
interior region where \(1+\lambda_Tg_t>0\). By Lemma~\ref{lem:part2-score-clt},
\[
    P_Tg_t=\Psi_T(\beta_0,\widehat\Lambda)=\Op(T^{-1/2}).
\]
By Lemma~\ref{lem:part2-variance-consistency},
\[
    P_Tg_t^2=\widehat V_T\to_p V>0.
\]
Expanding the KKT equation
\[
    0=P_T\left[\frac{g_t}{1+\lambda_Tg_t}\right]
\]
around \(\lambda_T=0\) gives
\[
    0=P_Tg_t-\lambda_TP_Tg_t^2
      +O_p\{\lambda_T^2P_T|g_t|^3\}.
\]
The third-moment bound in Assumption~\ref{ass:moment-control} implies the
remainder is \(O_p(\lambda_T^2)\). Since \(P_Tg_t=\Op(T^{-1/2})\) and
\(P_Tg_t^2\to_p V>0\), the root satisfies \(\lambda_T=\Op(T^{-1/2})\) and hence
\[
    \lambda_T=\frac{P_Tg_t}{P_Tg_t^2}+\op(T^{-1/2}).
\]
Combining \(\lambda_T=\Op(T^{-1/2})\) with the max-score bound gives
\[
    \max_{t\le T}|\lambda_Tg_t|
    \le |\lambda_T|\max_{t\le T}|g_t|=\op(1).
\]
Therefore the log EL ratio can be Taylor-expanded:
\[
\begin{aligned}
    \ell_T(\beta_0,\widehat\Lambda)
    &=2\sum_{t=1}^T\log(1+\lambda_Tg_t) \\
    &=2T\lambda_TP_Tg_t-T\lambda_T^2P_Tg_t^2
      +O_p\{T|\lambda_T|^3P_T|g_t|^3\}.
\end{aligned}
\]
The final term is \(\op(1)\), because \(\lambda_T=\Op(T^{-1/2})\) and
\(P_T|g_t|^3=\Op(1)\). Substituting
\(\lambda_T=P_Tg_t/P_Tg_t^2+\op(T^{-1/2})\) yields
\[
    \ell_T(\beta_0,\widehat\Lambda)
    =\frac{T(P_Tg_t)^2}{P_Tg_t^2}+\op(1)
    =\frac{T\{\Psi_T(\beta_0,\widehat\Lambda)\}^2}{\widehat V_T}+\op(1).
\]
\end{proof}

\subsection{Known-partition Wilks}

\begin{proposition}[Known-partition profile EL Wilks]
\label{prop:part2-known-partition-wilks}
Under the true partition \(\Lambda_0\), the profile empirical likelihood ratio
satisfies
\[
    \ell_T(\beta_0,\Lambda_0)\Rightarrow\chi_1^2.
\]
\end{proposition}

\begin{proof}
The true-partition versions of Lemmas~\ref{lem:part2-score-clt} and
\ref{lem:part2-variance-consistency} give
\[
    \sqrt T\Psi_T(\beta_0,\Lambda_0)\Rightarrow N(0,V),
    \qquad
    V_T^0\to_p V>0.
\]
The max-score and convex-hull parts of Assumptions~\ref{ass:moment-control} and
\ref{ass:convex-hull} apply to \(g_t(\beta_0,\Lambda_0)\), so the KKT expansion
from Part I gives
\[
    \ell_T(\beta_0,\Lambda_0)
    =\frac{T\{\Psi_T(\beta_0,\Lambda_0)\}^2}{V_T^0}+\op(1).
\]
Therefore
\[
    \ell_T(\beta_0,\Lambda_0)
    =\left\{
      \frac{\sqrt T\Psi_T(\beta_0,\Lambda_0)}{\sqrt{V_T^0}}
     \right\}^2+
     \op(1)
    \Rightarrow\chi_1^2.
\]
\end{proof}
\subsection{Part II consequences}

The uniform bridge gives the consistency input. If \(\widehat\beta_T\) is an
approximate sample root,
\[
    \Psi_T(\widehat\beta_T,\widehat\Lambda)=\op(1),
\]
then Lemma~\ref{lem:part2-uniform-bridge} gives
\[
    \Psi(\widehat\beta_T,\Lambda_0)=\op(1).
\]
Part I gives the unique population zero, so \(\widehat\beta_T\to_p\beta_0\).

For Wilks, Lemmas~\ref{lem:part2-score-clt},
\ref{lem:part2-variance-consistency}, and \ref{lem:part2-el-local-expansion}
give
\[
    \ell_T(\beta_0,\widehat\Lambda)
    =\frac{T\{\Psi_T(\beta_0,\widehat\Lambda)\}^2}{\widehat V_T}+\op(1),
    \qquad
    \frac{\sqrt T\Psi_T(\beta_0,\widehat\Lambda)}{\sqrt{\widehat V_T}}
    \Rightarrow N(0,1).
\]
Part III only merges these two displayed statements with the geometric KKT
reduction already proved in Part I.
\part{Merging Part}

\section{Merging Part: Moment Bridge, EL KKT, and \texorpdfstring{\(\chi_1^2\)}{chi-square}}

The merge step should not say that the unscaled moment bridge is itself
equivalent to Wilks. The correct statement is more precise:

\begin{center}
\fbox{\parbox{0.9\textwidth}{
The same projected moment \(\Psi_T\) is the interface for both targets.
Unscaled convergence of \(\Psi_T\) to \(\Psi\) merges with population uniqueness
to prove consistency. Scaled convergence of \(\Psi_T\), together with the EL KKT
quadratic reduction, merges to prove \(\chi_1^2\).}}
\end{center}

\subsection{Merge A: consistency}

Part I gives the population geometry
\[
    \Psi(\beta,\Lambda_0)
    =\innerP{M_0^PY_\beta^\circ}{v_0^P}
    =-(\beta-\beta_0)\Gamma
\]
under exogeneity and relevance. Hence
\[
    \Psi(\beta,\Lambda_0)=0
    \quad\Longleftrightarrow\quad
    \beta=\beta_0.
\]
Part II gives the stochastic bridge
\[
    \sup_{\beta\in\mathcal B}
    |\Psi_T(\beta,\widehat\Lambda)-\Psi(\beta,\Lambda_0)|=\op(1).
\]
Therefore any approximate sample root satisfies
\[
    \Psi_T(\widehat\beta_T,\widehat\Lambda)=\op(1)
    \quad\Longrightarrow\quad
    \widehat\beta_T\to_p\beta_0.
\]
This is the estimation target.

\subsection{Merge B: Wilks inference}

For inference, Part I gives the EL KKT representation
\[
    \ell_T(\beta,\widehat\Lambda)
    =2\sum_{t=1}^T\log\{1+\lambda_Tg_t(\beta,\widehat\Lambda)\},
\]
where \(\lambda_T\) solves
\[
    P_T\left[\frac{g_t(\beta,\widehat\Lambda)}
    {1+\lambda_Tg_t(\beta,\widehat\Lambda)}\right]=0.
\]
Part II verifies the local size and convex-hull conditions, so the KKT equation
can be expanded at \(\beta=\beta_0\):
\[
    \ell_T(\beta_0,\widehat\Lambda)
    =\frac{T\{\Psi_T(\beta_0,\widehat\Lambda)\}^2}
           {V_T(\widehat\Lambda)}+\op(1).
\]
Part II also gives
\[
    \frac{\sqrt T\Psi_T(\beta_0,\widehat\Lambda)}
         {\sqrt{V_T(\widehat\Lambda)}}
    \Rightarrow N(0,1).
\]
Combining the last two displays gives
\begin{equation}
    \ell_T(\beta_0,\widehat\Lambda)
    \Rightarrow \chi_1^2.
    \label{eq:target-wilks}
\end{equation}

Thus the ``equivalence'' is conditional on the geometric KKT reduction. Once
EL has been reduced to the square of the standardized projected moment, Wilks is
equivalent to the standardized score CLT. It is not equivalent to the unscaled
LLN bridge used for consistency.

\begin{center}
\fbox{\parbox{0.9\textwidth}{
\textbf{Proof-body split.} Geometry owns the model, nuisance tangent space,
projection KKT, residual score, population zero, and EL primal/dual KKT. Small
stochastic bounds only verify convergence and local size. The merging part then
turns the same projected moment into either consistency or Wilks.}}
\end{center}

\section{Optional Semiparametric Efficiency Theorem With Efficient Weight}

This optional block comes after the Wilks argument and is not part of the
minimum profile-sieve Wilks proof. The generic projected score with a residualized
instrument \(w\) gives valid orthogonal EL inference under the stochastic
conditions above. A Fisher-style semiparametric efficiency claim needs an extra
local likelihood model and the efficient residualized instrument.

The logic is: choose \(w=w^*\), prove efficient-score equivalence, then derive
the asymptotic linear expansion. The homoskedastic case reduces to \(w=X\); the
heteroskedastic oracle case uses \(w^*=X/\sigma^2(A)\).

Let
\[
    A=(J_0,W),
\]
where \(J_0\) is the true regime indicator. If the paper keeps the notation
\(A=(S,W)\), then every conditional projection \(E(\cdot\mid A)\) below should
be read as projection onto the true break-aware nuisance sigma-field generated
by the true regime and \(W\). This avoids treating the nuisance space as an
arbitrary continuous-time \(S\)-indexed space. The population model is
\[
    Y=\beta_0X+\mu_0(A)+U,
    \qquad
    \mu_0(A)\in\N_0^P=\N_{\Lambda_0}^P.
\]
For the efficiency calculation, assume the conditional Gaussian location model
\[
    U\mid X,A\sim N(0,\sigma^2(A)),
    \qquad
    0<c\le \sigma^2(A)\le C<\infty.
\]
The homoskedastic case is the special case \(\sigma^2(A)=\sigma^2\). The
population nuisance tangent directions generated by paths
\(\mu_t(A)=\mu_0(A)+t h(A)\), \(h\in\N_0^P\), are
\[
    \mathcal T_\mu
    =\overline{\left\{
      \frac{Uh(A)}{\sigma^2(A)}:h\in\N_0^P
     \right\}}^{L^2(P)}.
\]
The raw score for \(\beta\) is
\[
    S_\beta^0=\frac{UX}{\sigma^2(A)}.
\]

\begin{proposition}[Efficient score]
\label{prop:efficient-score}
Let
\[
    m_X(A)=E(X\mid A),
    \qquad
    \widetilde X=X-m_X(A).
\]
The efficient score for \(\beta\) is
\[
    S_{\mathrm{eff}}
    =\frac{U\widetilde X}{\sigma^2(A)},
\]
and the efficient information is
\[
    I_{\mathrm{eff}}
    =E(S_{\mathrm{eff}}^2)
    =E\left\{\frac{\widetilde X^2}{\sigma^2(A)}\right\}.
\]
\end{proposition}

\begin{proof}
The projection of \(S_\beta^0\) onto \(\mathcal T_\mu\) has the form
\(Uh^*(A)/\sigma^2(A)\). The function \(h^*\) minimizes
\[
    E\left[
      \left\{\frac{U}{\sigma^2(A)}\right\}^2
      \{X-h(A)\}^2
    \right].
\]
Since \(E(U^2\mid X,A)=\sigma^2(A)\), this criterion is
\[
    E\left[\frac{\{X-h(A)\}^2}{\sigma^2(A)}\right].
\]
Because \(\sigma^2(A)\) is \(A\)-measurable, conditional minimization gives
\(h^*(A)=E(X\mid A)=m_X(A)\). Therefore
\[
    \Pi(S_\beta^0\mid\mathcal T_\mu)
    =\frac{U m_X(A)}{\sigma^2(A)},
\]
and
\[
    S_{\mathrm{eff}}
    =S_\beta^0-\Pi(S_\beta^0\mid\mathcal T_\mu)
    =\frac{U\{X-m_X(A)\}}{\sigma^2(A)}.
\]
Finally,
\[
    E(S_{\mathrm{eff}}^2)
    =E\left\{\frac{U^2\widetilde X^2}{\sigma^4(A)}\right\}
    =E\left\{\frac{\widetilde X^2}{\sigma^2(A)}\right\},
\]
again using \(E(U^2\mid X,A)=\sigma^2(A)\).
\end{proof}

The efficient population weight is
\[
    w^*(O)=\frac{X}{\sigma^2(A)}.
\]
Since \(\sigma^2(A)\) is \(A\)-measurable,
\[
    E(w^*\mid A)=\frac{E(X\mid A)}{\sigma^2(A)}.
\]
Thus the population residualized efficient instrument is
\begin{equation}
    M_0^Pw^*
    =w^*-E(w^*\mid A)
    =\frac{\widetilde X}{\sigma^2(A)}.
    \label{eq:population-efficient-instrument}
\end{equation}
At \(\beta_0\), \(Y_{\beta_0}^\circ=Y-\beta_0X=\mu_0(A)+U\), and
\(M_0^P\mu_0=0\). Hence
\[
    M_0^PY_{\beta_0}^\circ\,M_0^Pw^*
    =U\frac{\widetilde X}{\sigma^2(A)}
    =S_{\mathrm{eff}}.
\]
This is the population reason the profile score becomes efficient when
\(w=w^*\). If \(\sigma^2(A)\) is constant, multiplying the moment by the
constant \(1/\sigma^2\) does not change the root, so \(w=X\) is the efficient
choice in the homoskedastic case.

For the sample statement, define
\[
    w_t^*=\frac{x_t}{\sigma^2(A_t)},
    \qquad
    \bm w^*=(w_1^*,\ldots,w_T^*)',
    \qquad
    \bm v_{K,T}^*=\bm M_0\bm w^*,
\]
where \(A_t=(J_{0,t},W_{t-1})\). Here \(J_{0,t}=j\) iff \(k_{j,0}<t\le k_{j+1,0}\), and \(\bm M_0=\bm M_{K,\Lambda_0}\). The efficient
sieve approximation target is
\begin{equation}
    \left\|\bm v_{K,T}^*
    -\left(\frac{\widetilde X_1}{\sigma^2(A_1)},\ldots,
            \frac{\widetilde X_T}{\sigma^2(A_T)}\right)'\right\|_T
    =\op(1).
    \label{eq:efficient-instrument-sieve-approximation}
\end{equation}
This follows from the same finite-sieve approximation and Gram-stability logic
used in Part II, provided the block sieve approximates
\(E(w^*\mid A)=E(X\mid A)/\sigma^2(A)\).

\begin{lemma}[Efficient-score equivalence]
\label{lem:efficient-score-equivalence}
Assume \eqref{eq:efficient-instrument-sieve-approximation} and the
true-partition efficient-weight bias condition
\[
    T^{-1/2}\bm r_{\mu,T}'\bm v_{K,T}^*=\op(1),
    \qquad
    \bm r_{\mu,T}=\bm M_0\bm\mu_0.
\]
Then
\[
    \sqrt T\Psi_T^*(\beta_0,\Lambda_0)
    =T^{-1/2}\sum_{t=1}^T S_{\mathrm{eff},t}+\op(1),
\]
where
\[
    \Psi_T^*(\beta,\Lambda)
    =\innerT{\bm M_{K,\Lambda}(\bm y-\beta\bm x)}
             {\bm M_{K,\Lambda}\bm w^*}.
\]
\end{lemma}

\begin{proof}
At \(\beta_0\),
\[
    \bm y_{\beta_0}^\circ=\bm\mu_0+\bm u.
\]
Therefore
\[
\begin{aligned}
    \sqrt T\Psi_T^*(\beta_0,\Lambda_0)
    &=T^{-1/2}(\bm M_0\bm y_{\beta_0}^\circ)'\bm v_{K,T}^* \\
    &=T^{-1/2}\bm u'\bm v_{K,T}^*
      +T^{-1/2}\bm r_{\mu,T}'\bm v_{K,T}^*.
\end{aligned}
\]
The second term is \(\op(1)\) by assumption. The first term equals
\[
\begin{aligned}
    T^{-1/2}\bm u'\bm v_{K,T}^*
    &=T^{-1/2}\sum_{t=1}^T U_t
      \frac{\widetilde X_t}{\sigma^2(A_t)}+\op(1) \\
    &=T^{-1/2}\sum_{t=1}^T S_{\mathrm{eff},t}+\op(1).
\end{aligned}
\]
by \eqref{eq:efficient-instrument-sieve-approximation} and the same
Cauchy-Schwarz/variance argument used for the score replacement in Part II.
\end{proof}

For fixed \(\Lambda\), the efficient-weight moment is linear in \(\beta\), with
\[
    \frac{\partial}{\partial\beta}\Psi_T^*(\beta,\Lambda)
    =-\innerT{\bm M_{K,\Lambda}\bm x}{\bm M_{K,\Lambda}\bm w^*}.
\]
Thus the derivative denominator in the root expansion is the empirical version
of the efficient information.

\begin{theorem}[Optional semiparametric efficiency with efficient weight]
\label{thm:semiparametric-efficiency}
Let \(\widehat\beta_T\) be a profile EL root, or an approximate profile moment
root, satisfying
\[
    \Psi_T^*(\widehat\beta_T,\widehat\Lambda)=\op(T^{-1/2}).
\]
Assume the efficient-score equivalence in
Lemma~\ref{lem:efficient-score-equivalence}, estimated-partition score
equivalence at the \(T^{-1/2}\) scale, and derivative equivalence:
\[
    \sqrt T\{\Psi_T^*(\beta_0,\widehat\Lambda)
      -\Psi_T^*(\beta_0,\Lambda_0)\}=\op(1),
\]
\[
    \innerT{\bm M_{K,\widehat\Lambda}\bm x}
            {\bm M_{K,\widehat\Lambda}\bm w^*}
    -\innerT{\bm M_0\bm x}{\bm v_{K,T}^*}=\op(1).
\]
Assume also
\[
    \innerT{\bm M_0\bm x}{\bm v_{K,T}^*}\to_p I_{\mathrm{eff}}>0.
\]
Then
\[
    \sqrt T(\widehat\beta_T-\beta_0)
    =I_{\mathrm{eff}}^{-1}T^{-1/2}\sum_{t=1}^T S_{\mathrm{eff},t}+\op(1),
\]
and consequently
\[
    \sqrt T(\widehat\beta_T-\beta_0)
    \Rightarrow N(0,I_{\mathrm{eff}}^{-1}).
\]
Thus the efficient-weight profile-sieve EL root attains the semiparametric
efficiency bound.
\end{theorem}

\begin{proof}
By a mean-value expansion of the projected moment around \(\beta_0\),
\[
\begin{aligned}
    0
    &=\Psi_T^*(\widehat\beta_T,\widehat\Lambda)+\op(T^{-1/2}) \\
    &=\Psi_T^*(\beta_0,\widehat\Lambda)
      -(\widehat\beta_T-\beta_0)
       \innerT{\bm M_{K,\widehat\Lambda}\bm x}
               {\bm M_{K,\widehat\Lambda}\bm w^*}
      +\op(T^{-1/2})+
       \op(|\widehat\beta_T-\beta_0|).
\end{aligned}
\]
The score-equivalence and derivative-equivalence assumptions replace the
estimated partition by the true partition. Hence
\[
    \sqrt T(\widehat\beta_T-\beta_0)
    =\left\{\innerT{\bm M_0\bm x}{\bm v_{K,T}^*}\right\}^{-1}
      \sqrt T\Psi_T^*(\beta_0,\Lambda_0)+\op(1).
\]
The denominator converges to
\[
    E\{(M_0^PX)(M_0^Pw^*)\}
    =E\left\{\widetilde X\frac{\widetilde X}{\sigma^2(A)}\right\}
    =I_{\mathrm{eff}}.
\]
Lemma~\ref{lem:efficient-score-equivalence} gives
\[
    \sqrt T\Psi_T^*(\beta_0,\Lambda_0)
    =T^{-1/2}\sum_{t=1}^T S_{\mathrm{eff},t}+\op(1).
\]
This proves the asymptotic linear representation. Since
\(E(S_{\mathrm{eff}}\mid X,A)=0\) and
\(E(S_{\mathrm{eff}}^2)=I_{\mathrm{eff}}\), the relevant martingale or iid CLT
gives
\[
    T^{-1/2}\sum_{t=1}^T S_{\mathrm{eff},t}
    \Rightarrow N(0,I_{\mathrm{eff}}).
\]
Therefore
\[
    \sqrt T(\widehat\beta_T-\beta_0)
    \Rightarrow N(0,I_{\mathrm{eff}}^{-1}).
\]
The standard Cauchy-Schwarz information bound gives
\(E(\phi^2)\ge I_{\mathrm{eff}}^{-1}\) for any regular influence function
\(\phi\) satisfying \(E\{\phi S_{\mathrm{eff}}\}=1\). The influence function
above is \(I_{\mathrm{eff}}^{-1}S_{\mathrm{eff}}\), whose variance is
\(I_{\mathrm{eff}}^{-1}\). It therefore attains the bound.
\end{proof}

\begin{remark}[Estimated variance plug-in]
The main theorem is cleanest with known \(\sigma^2(A)\), or with the
homoskedastic simplification \(w=X\). If \(\sigma^2(A)\) is estimated, define
\[
    \widehat w_t^*=\frac{x_t}{\widehat\sigma^2(A_t)},
    \qquad
    \widehat{\bm w}^*=(\widehat w_1^*,\ldots,\widehat w_T^*)'.
\]
Replacing \(\bm w^*\) by \(\widehat{\bm w}^*\) leaves the same asymptotic linear
representation if
\[
    \|\bm M_{K,\widehat\Lambda}(\widehat{\bm w}^*-\bm w^*)\|_T=\op(1),
\]
\[
    \sqrt T\left|\innerT{\bm M_{K,\widehat\Lambda}(\bm y-\beta_0\bm x)}
    {\bm M_{K,\widehat\Lambda}(\widehat{\bm w}^*-\bm w^*)}\right|=\op(1),
\]
and
\[
    \left|\innerT{\bm M_{K,\widehat\Lambda}\bm x}
    {\bm M_{K,\widehat\Lambda}(\widehat{\bm w}^*-\bm w^*)}\right|=\op(1).
\]
This plug-in statement is an appendix-level add-on; the main optional theorem
should use either \(w=X\) under homoskedasticity or the oracle heteroskedastic
weight \(w^*=X/\sigma^2(A)\).
\end{remark}
This theorem is conditional on the efficiency model and the efficient-weight
choice. With only martingale exogeneity and no conditional variance or local
likelihood structure, the workbook proves valid profile EL Wilks inference, but
not a Fisher-style semiparametric efficiency bound.

\end{document}
